\documentclass[11pt,a4paper]{article}

\usepackage[margin=1in]{geometry}
\usepackage{amsmath,amssymb,amsthm,mathtools}
\usepackage{array}
\usepackage{booktabs}
\usepackage{float}
\usepackage{graphicx}
\usepackage[hidelinks]{hyperref}
\usepackage{microtype}

\numberwithin{equation}{section}

\theoremstyle{plain}
\newtheorem{theorem}{Theorem}[section]
\newtheorem{lemma}[theorem]{Lemma}
\newtheorem{proposition}[theorem]{Proposition}
\newtheorem{corollary}[theorem]{Corollary}
\theoremstyle{definition}
\newtheorem{definition}[theorem]{Definition}
\theoremstyle{remark}
\newtheorem{remark}[theorem]{Remark}

\newcommand{\TT}{\mathbb{T}}
\newcommand{\Fold}{\mathrm{Fold}}
\newcommand{\leb}{\mathrm{Leb}}
\newcommand{\pref}{\mathrm{pref}}
\DeclareMathOperator{\TV}{TV}
\DeclareMathOperator{\KL}{KL}
\DeclareMathOperator{\supp}{supp}

\title{Zeckendorf folds, Sturmian rigidity, and Parry divergence on the golden-mean shift}
\author{Haobo Ma\thanks{AELF PTE LTD., \#14-02, Marina Bay Financial Centre Tower 1, 8 Marina Blvd, Singapore 018981, Singapore. \texttt{auric@aelf.io}.}
\and
Wenlin Zhang\thanks{National University of Singapore, Singapore. Corresponding author. \texttt{e1327962@u.nus.edu}.}}
\date{}

\begin{document}
\maketitle

\begin{abstract}
We study the canonical Zeckendorf fold $\Fold_m:\{0,1\}^m\to X_m$, where
$X_m$ is the length-$m$ language of the golden-mean shift. On the canonical
Sturmian slice $\beta=\alpha$, the restriction of $\Fold_m$ to realized
rotation factors is bijective and induces the classical higher-block
presentation of the Sturmian subshift. Using the endpoint collapse of
golden-mean Parry cylinders, we derive an exact formula for
$\KL(\nu\|\pi_m)$, sharp support-constrained minima for all R\'enyi orders,
and, for $m\ge 6$, the exact optimal injective placement of the realized
Sturmian law into $X_m$. For bounded-type slopes this yields
$$
D_q(\pi_m^{\mathrm{fold}}\|\pi_m)
=m\log\varphi-\log m+O_\alpha(1),
$$
uniformly in $q$. We further show that this finite-$m$ formula calibrates
the universal injective-placement barrier for lossless coarse-grainings into
primitive shifts of finite type, and that the Fibonacci weighting is uniquely
characterized by a natural contiguous numeration axiom. A short final section
records the conjugacy-invariant transfer of the construction to irrational
circle diffeomorphisms and includes one numerical illustration.
\end{abstract}

\medskip
\noindent\textbf{Keywords.} Sturmian subshift, Zeckendorf numeration,
golden-mean shift, Parry measure, R\'enyi divergence, shifts of finite type

\medskip
\noindent\textbf{2020 Mathematics Subject Classification.}
37B10 (primary), 37A55, 11B39, 94A17

\section{Introduction}\label{sec:introduction}

This paper studies the canonical Zeckendorf fold, a deterministic map
$\Fold_m:\{0,1\}^m\to X_m$ that sends binary words to the length-$m$
language $X_m$ of the golden-mean shift by Fibonacci-weighted evaluation and
Zeckendorf normalization. Three structural phenomena are proved.

First, on the Sturmian slice $\beta=\alpha$, the fold is bijective on
realized rotation factors and the induced sequence map is topologically
conjugate to the classical higher-block presentation of the Sturmian
subshift. Minimality, unique ergodicity, zero entropy, and linear word
complexity are therefore preserved exactly.

Second, the endpoint collapse of golden-mean Parry cylinders yields a closed
formula for $\KL(\nu\|\pi_m)$ and, more generally, sharp
support-constrained Parry divergence minima for all R\'enyi orders. The
formula depends only on the endpoints of the admissible words, and its
universality over $q\in[0,\infty]$ is a consequence of the scalar
endpoint statistic peculiar to the golden-mean transition matrix. For
$m\ge 6$, the exact optimal injective placement of the realized Sturmian
law into $X_m$ is determined: the infimum over all injections
$\iota:S_m(\alpha)\hookrightarrow X_m$ equals
\[
L_q(m)=(m-1)\log\varphi-\log\pi(0)-H_q(\mu_m),
\]
and every injection whose image lies in the heaviest endpoint class
$X_m^{00}$ attains it. Moreover, the full $O(1)$ defect between the
canonical folded law and the injective optimum is shown to be an exact
escort-pressure of a boundary cocycle
(Theorem~\ref{thm:full-renyi-defect-formula}), not an entropy residual.

Third, for bounded-type slopes the folded law satisfies
\[
D_q(\pi_m^{\mathrm{fold}}\|\pi_m)
=m\log\varphi-\log m+O_\alpha(1),
\qquad q\in[0,\infty],
\]
and the canonical fold remains within a bounded additive distance of the
injective optimum. This finite-$m$ formula calibrates the universal
injective-placement barrier for primitive shifts of finite type
(Theorem~\ref{thm:primitive-sft-injective-barrier}): any deterministic map
that is injective on the realized source retains the order-uniform
divergence rate $D_q^{\mathrm{opt}}=mh-H_q(\mu_m)+O(1)$. The golden-mean
benchmark is exceptional because this universal envelope becomes finite-$m$
exact, the optimizer is explicit, and the actual canonical fold remains
within a bounded additive distance of that optimizer. Moreover, the
Fibonacci weighting is proved to be the unique choice satisfying a natural
contiguous numeration axiom (Theorem~\ref{thm:fibonacci-weight-uniqueness}).

The ingredients are individually classical. Sturmian codings, factor
complexity, and factor frequencies go back to
\cite{MorseHedlund1940,Slater1950DistributionIntegers,Slater1967GapsSteps}
and fit naturally into the substitution and symbolic-dynamical framework
\cite{LindMarcus1995,QueffelecSubstitution,Fogg2002}. The Parry measure and
the thermodynamic formalism of shifts of finite type are standard
\cite{LindMarcus1995,WaltersErgodicTheory,ParryPollicott1990Zeta}. Fibonacci
numeration and finite-state normalization are likewise classical
\cite{Zeckendorf1972,Frougny1991FibonacciFiniteAutomata,Frougny1992NumbersAutomata},
with related Sturmian and Ostrowski numeration viewpoints in
\cite{Berthe2001Ostrowski,BaranwalaSchaefferShallit2021OstrowskiAutomatic}.
What is specific to the present construction is that these ingredients
combine cleanly: the fold is dynamically lossless on the realized Sturmian
source, while the comparison with the Parry measure is explicit enough to
yield exact finite-$m$ optimization statements. On the numeration side,
Ostrowski and related continued-fraction systems
\cite{Berthe2001Ostrowski,BaranwalaSchaefferShallit2021OstrowskiAutomatic}
provide arithmetic codings with rich dynamical content; the present paper
uses the golden branch of that numeration-theoretic viewpoint as an exactly
analyzable coarse-graining whose information geometry relative to the Parry
law can be solved sharply.

For primitive one-step shifts of finite type, support-constrained R\'enyi
optimization remains finite-$m$ exact at the endpoint-class level
(Theorem~\ref{thm:primitive-sft-endpoint-envelope}). The golden-mean case
is exceptional only because its endpoint spectrum collapses to three
Fibonacci classes, so the general endpoint envelope closes in elementary
form. This upgrades the role of the golden-mean setting from a single
benchmark to the most explicit instance of a general endpoint-class theory
for primitive shifts. By topological conjugacy, the entire construction
transfers verbatim to nonlinear circle diffeomorphisms with irrational
rotation number.

The paper is organized as follows. Section~\ref{sec:fold-language} introduces
the canonical Zeckendorf fold, the golden-mean language, and the Parry
reference measure. Section~\ref{sec:sturmian-rigidity} proves the
realized-support bijection and the higher-block conjugacy on the Sturmian
source. Section~\ref{sec:parry-divergence} derives the exact Parry divergence
identity, the sharp support barriers, and the endpoint-class optimization.
Section~\ref{sec:injective-placement} establishes the optimal injective
relabeling, the bounded-type asymptotics, and the primitive-SFT extension.
Section~\ref{sec:conjugacy} records the conjugacy-invariant transfer to circle
diffeomorphisms and one numerical illustration. Section~\ref{sec:concluding}
collects concluding remarks. Three appendices supply auxiliary discrepancy and
empirical entropy bounds (Appendix~\ref{app:discrepancy}), fiber-uniform lifts
and I-projections (Appendix~\ref{app:fiber-lift}), and the local Zeckendorf
normalization together with the uniqueness of the Fibonacci weighting
(Appendix~\ref{app:canonical-fold}).

\section{The Canonical Zeckendorf Fold and the Golden-Mean Language}
\label{sec:fold-language}

\subsection{Rotations, Windows, and the Folded Observable}

Let $\alpha\in(0,1)$ be irrational, let $\beta\in(0,1)$, and let
$x_0\in\TT:=\mathbb{R}/\mathbb{Z}$.  Define the Kronecker orbit
\[
x_t:=x_0+t\alpha \pmod 1,\qquad t\ge 0,
\]
and the binary readout
\[
s_t:=\mathbf{1}_{[0,\beta)}(x_t)\in\{0,1\}.
\]
For a fixed resolution $m\ge 1$, let
\[
\Omega_m:=\{0,1\}^m,\qquad
X_m:=\{u=(u_1,\dots,u_m)\in\{0,1\}^m:\ u_i u_{i+1}=0\text{ for }1\le i<m\}.
\]

\begin{definition}[Canonical fold]\label{def:canonical-fold}
Let \(F_1=F_2=1\) and \(F_{k+2}=F_{k+1}+F_k\). For
\(\omega=\omega_1\cdots\omega_m\in\Omega_m\), define its Fibonacci value by
\[
N(\omega):=\sum_{k=1}^m \omega_k F_{k+1}.
\]
By Zeckendorf's theorem \cite{Zeckendorf1972}, there is a unique sequence
\(Z(N(\omega))=(c_1,c_2,\dots)\in\{0,1\}^{\mathbb{N}}\) such that
\[
N(\omega)=\sum_{k\ge 1} c_k F_{k+1},
\qquad
c_kc_{k+1}=0\ \text{ for all }k\ge 1.
\]
Let
\[
\pref_m(c_1,c_2,\dots):=(c_1,\dots,c_m)\in X_m.
\]
The canonical fold at resolution \(m\) is
\[
\Fold_m(\omega):=\pref_m(Z(N(\omega)))\in X_m.
\]
\end{definition}

\begin{proposition}[Basic properties of the canonical fold]\label{prop:canonical-fold-basic}
For every \(m\ge 1\), the map \(\Fold_m:\Omega_m\to X_m\) is well defined,
surjective, and satisfies \(\Fold_m|_{X_m}=\mathrm{id}\). Moreover,
\(\Fold_m\) admits an order-independent local implementation by
value-preserving Zeckendorf rewrites.
\end{proposition}

\begin{proof}
Well-definedness is immediate from Definition~\ref{def:canonical-fold}. If
\(\omega\in X_m\), then \(\omega\) is already a legal Zeckendorf prefix, so the
uniqueness of the Zeckendorf expansion gives \(\Fold_m(\omega)=\omega\). Hence
\(\Fold_m\) fixes \(X_m\) pointwise, and surjectivity follows by taking
\(\omega=x\in X_m\) for any \(x\in X_m\). The order-independent local
implementation is recorded in Appendix~\ref{app:canonical-fold}.
\end{proof}

\begin{proposition}[Residue characterization of the canonical fold]\label{prop:fold-residue-characterization}
For every \(m\ge 1\) and \(\omega\in\Omega_m\), write
\[
N(\omega)=qF_{m+2}+r,
\qquad q\in\{0,1\},
\qquad 0\le r<F_{m+2}.
\]
Then
\[
\Fold_m(\omega)=\pref_m(Z(r)).
\]
Equivalently, for \(x\in X_m\),
\[
\Fold_m(\omega)=x
\iff
N(\omega)\equiv N(x)\pmod{F_{m+2}}.
\]
\end{proposition}

\begin{proof}
Since
\[
N(\omega)\le \sum_{k=1}^m F_{k+1}=F_{m+3}-2<F_{m+2}+F_{m+2},
\]
the quotient \(q\) is either \(0\) or \(1\). If \(q=0\), then \(r=N(\omega)\)
and the claim is immediate.

Assume \(q=1\). Then
\[
0\le r=N(\omega)-F_{m+2}\le F_{m+1}-2<F_{m+1},
\]
so the Zeckendorf expansion \(Z(r)\) has no nonzero digit in position \(m\) or
beyond. Hence appending a \(1\) in position \(m+1\) produces another legal
Zeckendorf digit string, namely the expansion of \(N(\omega)=F_{m+2}+r\). Its
first \(m\) digits are therefore exactly those of \(Z(r)\), which proves the
first statement.

If \(x\in X_m\), then
\[
N(x)\le \sum_{j=0}^{\lfloor (m-1)/2\rfloor}F_{m-2j+1}<F_{m+2},
\]
so \(x=\pref_m(Z(N(x)))\). Applying the first statement with
\(r\equiv N(\omega)\pmod{F_{m+2}}\) yields
\[
\Fold_m(\omega)=x
\iff
\pref_m(Z(r))=\pref_m(Z(N(x)))
\iff
r=N(x),
\]
which is equivalent to the congruence condition.
\end{proof}

\begin{definition}[Microstates and folded histograms]
For $t\ge 0$, the length-$m$ microstate is
\[
\omega_m(t):=s_t s_{t+1}\cdots s_{t+m-1}\in\Omega_m.
\]
The folded output is
\[
y_m(t):=\Fold_m(\omega_m(t))\in X_m.
\]
For a sample length $N$, define the empirical microstate law and folded
histogram by
\[
\hat\mu_{m,N}(u):=\frac{1}{N}\sum_{t=0}^{N-1}\mathbf{1}_{\{\omega_m(t)=u\}},
\qquad u\in\Omega_m,
\]
\[
\hat\pi_{m,N}(v):=\frac{1}{N}\sum_{t=0}^{N-1}\mathbf{1}_{\{y_m(t)=v\}},
\qquad v\in X_m.
\]
\end{definition}

Because irrational rotations are uniquely ergodic, the empirical microstate
law converges to a deterministic target law induced by Lebesgue measure.

\begin{definition}[Target laws]
Let $\mu_m$ be the length-$m$ microstate distribution induced by Lebesgue
measure on $\TT$:
\[
\mu_m(u):=\leb\{x\in\TT:\omega_m(x;0)=u\},\qquad u\in\Omega_m.
\]
Its folded pushforward is
\[
\pi_m^{\mathrm{fold}}:=(\Fold_m)_*\mu_m\in\Delta(X_m).
\]
\end{definition}

\begin{remark}[Spaces, maps, and laws]\label{rem:model-spaces-maps}
The basic objects may be summarized by
\[
\begin{array}{ccc}
\Omega_m & \xrightarrow{\Fold_m} & X_m\\[0.3em]
\mu_m \downarrow\phantom{(\Fold_m)_*} &  & \phantom{(\Fold_m)_*}\downarrow \pi_m^{\mathrm{fold}}\\[0.3em]
\hat\mu_{m,N} & \xrightarrow{(\Fold_m)_*} & \hat\pi_{m,N}
\end{array}
\qquad
\begin{array}{ccc}
\Omega_{m+1} & \xrightarrow{\Fold_{m+1}} & X_{m+1}\\[0.3em]
\pref_m \downarrow &  & \downarrow \rho_{m+1\to m}\\[0.3em]
\Omega_m & \xrightarrow{\Fold_m} & X_m.
\end{array}
\]
The left diagram separates microstate space \(\Omega_m\) from folded space
\(X_m\). The right diagram records the truncation interface whose failure to
commute is measured by the projective residual in
Appendix~\ref{app:discrepancy}.
\end{remark}

\begin{remark}[Scope]\label{rem:scope-fold-assumptions}
We fix the canonical fold throughout. The Sturmian rigidity and higher-block
conjugacy of Section~\ref{sec:sturmian-rigidity} use the Sturmian slice
$\beta=\alpha$. The sharp Parry divergence identity of
Section~\ref{sec:parry-divergence} is a statement about the Parry reference
family on $X_m$ and is independent of the fold once a law on $X_m$ has been
specified. The injective placement results of
Section~\ref{sec:injective-placement} combine the two: the low-complexity
support created by the folded rotation coding meets the ambient equilibrium
structure of the golden-mean shift.
\end{remark}

\subsection{The Parry Measure}

As a reference measure on \(X_m\) we use the Parry measure of the golden-mean
shift. Let
\[
A=\begin{pmatrix}1&1\\1&0\end{pmatrix},\qquad
\varphi=\frac{1+\sqrt5}{2}.
\]
The associated two-state Parry chain has transition matrix
\[
P=\begin{pmatrix}\varphi^{-1} & \varphi^{-2}\\ 1&0\end{pmatrix},
\]
and stationary distribution
\[
\pi(1)=\frac{1}{\varphi^2+1},\qquad
\pi(0)=1-\pi(1).
\]
For $u=u_1\cdots u_m\in X_m$, define the length-$m$ Parry cylinder law
\[
\pi_m(u):=\pi(u_1)\prod_{j=1}^{m-1}P_{u_j u_{j+1}}.
\]
We emphasize that \(\pi_m\) is used as a reference measure only. We do not
assume that \(\pi_m^{\mathrm{fold}}=\pi_m\) for irrational rotations.
Since \(|X_m|=F_{m+2}\), the ambient shift has topological entropy
\[
h_{\mathrm{top}}(X)=\lim_{m\to\infty}\frac1m\log |X_m|=\log\varphi,
\]
and the Parry measure is its unique measure of maximal entropy
\cite{LindMarcus1995,WaltersErgodicTheory}.

\begin{proposition}[Endpoint collapse of Parry cylinders]\label{prop:parry-endpoint-collapse}
For every legal word \(u=u_1\cdots u_m\in X_m\),
\[
\pi_m(u)=
\begin{cases}
\pi(0)\varphi^{-(m-1)}, & (u_1,u_m)=(0,0),\\
\pi(0)\varphi^{-m}, & (u_1,u_m)\in\{(0,1),(1,0)\},\\
\pi(0)\varphi^{-(m+1)}, & (u_1,u_m)=(1,1).
\end{cases}
\]
Equivalently,
\[
\pi_m(u)
=
\pi(0)\varphi^{-m+\mathbf{1}_{\{u_1=0\}}-\mathbf{1}_{\{u_m=1\}}}.
\]
In particular, the Parry cylinder weight depends only on the endpoints of the
word.
\end{proposition}

\begin{proof}
Let \(r_0:=\varphi\) and \(r_1:=1\). For every legal transition
\((i,j)\in\{(0,0),(0,1),(1,0)\}\), one has
\[
P_{ij}=\frac{r_j}{\varphi r_i}.
\]
Hence for a legal word \(u=u_1\cdots u_m\),
\[
\prod_{j=1}^{m-1}P_{u_j u_{j+1}}
=
\varphi^{-(m-1)}\frac{r_{u_m}}{r_{u_1}}.
\]
Multiplying by \(\pi(u_1)\) and using \(\pi(1)=\pi(0)\varphi^{-2}\) gives the
three cases above.
\end{proof}

\subsection{Sturmian Cylinder Bounds}

The following continued-fraction bounds on Sturmian cylinder lengths are used
in Sections~\ref{sec:sturmian-rigidity}
and~\ref{sec:injective-placement}.

\begin{theorem}[Bounds for Sturmian cylinder lengths]\label{thm:spg-sturmian-cylinder-measure-continued-fraction-sandwich}
Let $\alpha=[0;a_1,a_2,\dots]$ and let $p_n/q_n$ be its convergents.  If
$q_n\le t<q_{n+1}$, then every length-$t$ Sturmian cylinder $C(w)$ satisfies
\[
\frac{1}{q_{n+1}+q_n}
\;<\;
\leb(C(w))
\;<\;
\frac{1}{q_n}+\frac{1}{q_{n+1}}
\;<\;
\frac{2}{q_n}.
\]
Equivalently,
\[
\log\!\Bigl(\frac{q_n}{2}\Bigr)
\;<\;
-\log\leb(C(w))
\;<\;
\log(q_{n+1}+q_n).
\]
If $\alpha$ has bounded partial quotients, then
$-\log\leb(C(w))=\log t+O(1)$ uniformly in $w$.
\end{theorem}

\begin{proof}
The endpoints of length-$t$ cylinders are among the orbit points
$\{-k\alpha:0\le k\le t\}$, so cylinder lengths are the adjacent gaps in this
finite rotation set.  Standard continued-fraction gap estimates and the
three-gap theorem give
\[
\min_{1\le k\le t}\|k\alpha\|
=
\|q_n\alpha\|
\in
\left(\frac{1}{q_{n+1}+q_n},\frac{1}{q_{n+1}}\right),
\]
while the largest gap is at most
$\|q_{n-1}\alpha\|+\|q_n\alpha\|<q_n^{-1}+q_{n+1}^{-1}$; see
\cite{Slater1950DistributionIntegers,Slater1967GapsSteps}.  Taking negative
logarithms yields the second display.
\end{proof}

\section{Sturmian Rigidity and Higher-Block Conjugacy}
\label{sec:sturmian-rigidity}

In the canonical Sturmian slice, the fold does not merge the words that are
actually realized by the rotation. The reason is arithmetic: along the
Sturmian cylinder partition, the Fibonacci values of the factors remain
distinct modulo \(F_{m+2}\).

\begin{theorem}[Sturmian no-collision rigidity]\label{thm:sturmian-no-collision-rigidity}
Assume \(\beta=\alpha\), and define
\[
S_m(\alpha):=\supp(\mu_m)\subset\Omega_m,
\qquad
Y_m:=\Fold_m(S_m(\alpha))\subset X_m.
\]
Then the restriction
\[
\Fold_m\!\upharpoonright_{S_m(\alpha)}:S_m(\alpha)\to Y_m
\]
is a bijection.

More precisely, let
\[
\Xi_m:=\{(1-j)\alpha \bmod 1:\ 0\le j\le m\}\subset\TT,
\]
let \(I_0,\dots,I_m\) be the open atoms of \(\TT\setminus \Xi_m\) in
counterclockwise order, and choose \(x_\ell\in I_\ell\). Writing
\[
u^{(\ell)}:=\omega_m(x_\ell;0)\in S_m(\alpha),
\qquad
V_\ell:=N(u^{(\ell)}),
\]
the cyclic jump multiset
\[
\Delta_\ell:=V_{\ell+1}-V_\ell
\qquad (\ell\!\!\!\!\pmod{m+1})
\]
consists of exactly one positive jump \(+F_{m+1}\) and the \(m\) negative
jumps
\[
-1,\ -F_1,\ -F_2,\ \dots,\ -F_{m-1},
\]
counted with multiplicity. After a cyclic relabeling that starts immediately to
the right of the unique positive jump, the values \(V_\ell\) form a strictly
decreasing integer chain. In particular, the residues \(V_\ell\bmod F_{m+2}\)
are pairwise distinct.
\end{theorem}

\begin{proof}
The proof has two steps: compute the Fibonacci jump produced by crossing each
Sturmian boundary atom, and then observe that after the unique positive jump
the remaining cycle is a strictly decreasing chain whose total drop is smaller
than \(F_{m+2}\), so residue collisions cannot occur. The boundary bookkeeping
is the same one that underlies the standard factor partition of Sturmian words;
see \cite{LindMarcus1995,Fogg2002,Berthe2001Ostrowski}.

For \(0\le k\le m-1\), the \(k\)-th symbol of \(\omega_m(x;0)\) is
\[
\mathbf{1}_{[0,\alpha)}(x+k\alpha),
\]
so its discontinuities occur at the two points \(-k\alpha\) and
\(\alpha-k\alpha=(1-k)\alpha\) on \(\TT\). The union of all such discontinuity
points is exactly \(\Xi_m\). Hence \(\omega_m(x;0)\) is constant on each atom
\(I_\ell\), and \(S_m(\alpha)=\{u^{(\ell)}:0\le \ell\le m\}\).

Fix a boundary point \(\xi_j:=(1-j)\alpha\in\Xi_m\).
If \(1\le j\le m-1\), then crossing \(\xi_j\) from left to right changes the
symbols in positions \(j-1\) and \(j\) from \(01\) to \(10\), while all other
positions stay fixed: among the translated points \(x+k\alpha\), only the two
with \(k=j-1\) and \(k=j\) meet the coding boundary at the instant \(x=\xi_j\).
Therefore the Fibonacci value jump is
\[
F_{j+1}-F_{j+2}=-F_j.
\]
At \(\xi_0=\alpha\), only the first symbol changes, from \(1\) to \(0\), so
the jump is \(-F_1=-1\). At \(\xi_m=(1-m)\alpha\), only the last symbol
changes, from \(0\) to \(1\), so the jump is \(+F_{m+1}\). This proves the
stated jump multiset.

Now start with the atom immediately to the right of \(\xi_m\). Along the
remaining cycle, every subsequent jump is negative, so the corresponding values
form a strictly decreasing integer chain. Since
\[
1+\sum_{j=1}^{m-1}F_j=F_{m+1},
\]
the total drop before the cycle closes is exactly \(F_{m+1}<F_{m+2}\). Hence
no two values in this chain can be congruent modulo \(F_{m+2}\).
Proposition~\ref{prop:fold-residue-characterization} shows that
\(\Fold_m\) is constant on residue classes modulo \(F_{m+2}\), so its
restriction to \(S_m(\alpha)\) is injective. Bijectivity onto \(Y_m\) is then
tautological.
\end{proof}

\begin{remark}[What the rigidity proof uses]
\label{rem:sturmian-rigidity-special}
The proof uses three ingredients that are all special to the canonical slice
\(\beta=\alpha\): the Sturmian partition of the torus into exactly \(m+1\)
atoms, the explicit Fibonacci weighting in \(N(\cdot)\), and the residue
characterization of the canonical fold modulo \(F_{m+2}\). The argument is
therefore more rigid than a generic coding proof, which is why an analogous
jumps-and-residues theorem is not claimed for non-canonical folds.
\end{remark}

\begin{corollary}[Exact relabeling on the realized Sturmian language]\label{cor:sturmian-exact-relabeling}
Assume \(\beta=\alpha\), and let
\[
B_m:=\Fold_m\!\upharpoonright_{S_m(\alpha)}:S_m(\alpha)\to Y_m.
\]
Then \(B_m\) is a bijection, and for every \(N\ge 1\),
\[
\pi_m^{\mathrm{fold}}=(B_m)_*\mu_m,
\qquad
\hat\pi_{m,N}=(B_m)_*\hat\mu_{m,N}.
\]
Consequently,
\[
\TV(\hat\pi_{m,N},\pi_m^{\mathrm{fold}})
=
\TV(\hat\mu_{m,N},\mu_m),
\]
\[
\KL(\hat\pi_{m,N}\|\pi_m^{\mathrm{fold}})
=
\KL(\hat\mu_{m,N}\|\mu_m),
\]
and, writing \(H_q\) for the R\'enyi entropy of order \(q\), for every
\(q\in[0,\infty]\),
\[
H_q(\hat\pi_{m,N})=H_q(\hat\mu_{m,N}),
\qquad
H_q(\pi_m^{\mathrm{fold}})=H_q(\mu_m).
\]
\end{corollary}

\begin{proof}
The pushforward identities follow from Theorem~\ref{thm:sturmian-no-collision-rigidity}.
Since \(B_m\) is a bijection between the realized supports, total variation,
relative entropy, and all R\'enyi entropies are invariant under the relabeling
\(B_m\).
\end{proof}

\subsection{Higher-Block Conjugacy on the Sturmian Source}

\begin{theorem}[Higher-block conjugacy on the Sturmian source]
\label{thm:sequence-level-transparency}
Assume \(\beta=\alpha\), let \(X_\alpha\subset\{0,1\}^{\mathbb{Z}}\) be the
Sturmian subshift generated by the irrational rotation coding, and fix
\(m\ge 1\). Define
\[
(\Phi_m(s))_t:=\Fold_m(s_t s_{t+1}\cdots s_{t+m-1}),
\qquad t\in\mathbb{Z}.
\]
Then \(\Phi_m:X_\alpha\to \mathcal{Y}_m:=\Phi_m(X_\alpha)\) is a topological
conjugacy. Equivalently, \(\mathcal{Y}_m\) is a letter relabeling of the
classical \(m\)-block presentation of \(X_\alpha\).

In particular, \(\mathcal{Y}_m\) is minimal and uniquely ergodic,
\[
h_{\mathrm{top}}(\mathcal{Y}_m)=0,
\]
and, writing \(p_Z(n)\) for the word complexity of a subshift \(Z\), one has
\[
p_{\mathcal{Y}_m}(n)=p_{X_\alpha}(n+m-1)=n+m
\qquad (n\ge 1).
\]
\end{theorem}

\begin{proof}
Let \(\Psi_m:X_\alpha\to S_m(\alpha)^{\mathbb{Z}}\) be the classical higher-block
code
\[
(\Psi_m(s))_t:=s_t s_{t+1}\cdots s_{t+m-1}.
\]
Standard higher-block theory gives that \(\Psi_m\) is a topological conjugacy
from \(X_\alpha\) onto its \(m\)-block presentation
\cite{LindMarcus1995}. By Corollary~\ref{cor:sturmian-exact-relabeling},
\[
B_m:=\Fold_m\!\upharpoonright_{S_m(\alpha)}:S_m(\alpha)\to Y_m
\]
is a bijection. Hence the coordinatewise relabeling
\[
(\widetilde B_m(z))_t:=B_m(z_t),
\qquad z\in \Psi_m(X_\alpha),
\]
is a topological conjugacy from the \(m\)-block presentation onto
\(\mathcal{Y}_m\).
Because
\[
\Phi_m=\widetilde B_m\circ \Psi_m,
\]
\(\Phi_m\) is itself a topological conjugacy.

Minimality and unique ergodicity are preserved under conjugacy, and the
irrational rotation coding is classically minimal and uniquely ergodic. The
same is therefore true of \(\mathcal{Y}_m\). Likewise, topological entropy is
preserved, so \(h_{\mathrm{top}}(\mathcal{Y}_m)=h_{\mathrm{top}}(X_\alpha)=0\).

A word of length \(n\) in the \(m\)-block presentation corresponds to a word of
length \(n+m-1\) in \(X_\alpha\), and the relabeling \(\widetilde B_m\) does
not change word counts. Therefore
\[
p_{\mathcal{Y}_m}(n)=p_{X_\alpha}(n+m-1).
\]
Since \(X_\alpha\) is Sturmian, \(p_{X_\alpha}(k)=k+1\) for every \(k\ge 1\),
which gives \(p_{\mathcal{Y}_m}(n)=n+m\).
\end{proof}

\begin{corollary}[Pure degeneracy ledger on the Sturmian slice]\label{cor:sturmian-pure-degeneracy-ledger}
Assume \(\beta=\alpha\), and let \(\mu_m^e:=Q_m\pi_m^{\mathrm{fold}}\). Then
\[
\KL(\mu_m\|\mu_m^e)
=
\mathbb{E}_{\pi_m^{\mathrm{fold}}}\bigl[\log d_m(Y)\bigr].
\]
Likewise, for every \(N\ge 1\), if \(\hat\mu_{m,N}^e:=Q_m\hat\pi_{m,N}\), then
\[
\KL(\hat\mu_{m,N}\|\hat\mu_{m,N}^e)
=
\mathbb{E}_{\hat\pi_{m,N}}\bigl[\log d_m(Y)\bigr].
\]
\end{corollary}

\begin{proof}
Corollary~\ref{cor:sturmian-exact-relabeling} gives
\[
H(\mu_m)=H(\pi_m^{\mathrm{fold}}),
\qquad
H(\hat\mu_{m,N})=H(\hat\pi_{m,N}).
\]
Applying Proposition~\ref{prop:fold-entropy-ledger} to \(\mu_m\) and
\(\hat\mu_{m,N}\) yields the two identities.
\end{proof}

A complementary microstate lift formalism is collected in
Appendix~\ref{app:fiber-lift}. It records the fiber-uniform lift family, the
exact entropy ledger used in Corollary~\ref{cor:sturmian-pure-degeneracy-ledger},
and an optional Pythagorean identity for relative entropy against lifted macro
models.

\section{Exact Parry Divergence and Sharp Support Barriers}
\label{sec:parry-divergence}

The Parry measure is the unique maximal-entropy equilibrium state of the
ambient golden-mean shift, whose topological entropy is
\(h_{\mathrm{top}}(X)=\log\varphi\). This section quantifies how far
folded rotation laws sit from ambient equilibrium when their effective support
is much smaller than the full shift language. The exact identities below rely
on endpoint collapse for the golden-mean Parry cylinders, and no comparable
closed formula is claimed for arbitrary shifts of finite type.

\begin{definition}[R\'enyi entropy and divergence]\label{def:renyi-entropy-divergence}
Let \(E\) be a finite set. For a probability law \(\nu\) on \(E\), define
\[
H_q(\nu):=
\begin{cases}
\log |\supp(\nu)|,
& q=0,\\[0.4em]
\dfrac{1}{1-q}\log\!\Bigl(\sum_{x\in E}\nu(x)^q\Bigr),
& q\in(0,\infty)\setminus\{1\},\\[0.8em]
-\sum_{x\in E}\nu(x)\log \nu(x),
& q=1,\\[0.4em]
-\log \max_{x\in E}\nu(x),
& q=\infty.
\end{cases}
\]
For probability laws \(\nu\) and \(\rho\) on \(E\), define
\[
D_q(\nu\|\rho):=
\begin{cases}
-\log \rho(\supp(\nu)),
& q=0,\\[0.4em]
\dfrac{1}{q-1}\log\!\Bigl(\sum_{x\in E}\nu(x)^q\rho(x)^{1-q}\Bigr),
& q\in(0,\infty)\setminus\{1\},\\[0.8em]
\KL(\nu\|\rho),
& q=1,\\[0.4em]
\log \max_{x:\,\nu(x)>0}\dfrac{\nu(x)}{\rho(x)},
& q=\infty,
\end{cases}
\]
with the convention \(0\log 0:=0\), and with \(D_q(\nu\|\rho):=\infty\) for
\(q>0\) if \(\nu\) is not absolutely continuous with respect to \(\rho\).
\end{definition}

\subsection{Exact Parry Identity}

\begin{theorem}[Exact divergence identity relative to the Parry measure]\label{thm:parry-kl-identity}
Let \(\nu\) be a probability measure on \(X_m\), and define the endpoint masses
\[
\eta_{ij}(\nu):=\nu\bigl(\{u\in X_m:u_1=i,\ u_m=j\}\bigr),
\qquad i,j\in\{0,1\}.
\]
If
\[
H(\nu):=-\sum_{u\in X_m}\nu(u)\log \nu(u)
\]
denotes the Shannon entropy of \(\nu\), then
\[
\KL(\nu\|\pi_m)
=
m\log\varphi-H(\nu)-\log\pi(0)
{}+(\eta_{11}(\nu)-\eta_{00}(\nu))\log\varphi.
\]
Consequently,
\[
m\log\varphi-H(\nu)-\log\pi(0)-\log\varphi
\le
\KL(\nu\|\pi_m)
\le
m\log\varphi-H(\nu)-\log\pi(0)+\log\varphi.
\]
\end{theorem}

\begin{proof}
By Proposition~\ref{prop:parry-endpoint-collapse},
\[
\log \pi_m(u)
=
\log\pi(0)-m\log\varphi
{}+\bigl(\mathbf{1}_{\{u_1=0\}}-\mathbf{1}_{\{u_m=1\}}\bigr)\log\varphi.
\]
Therefore
\[
\begin{aligned}
\KL(\nu\|\pi_m)
&=
\sum_{u\in X_m}\nu(u)\log\frac{\nu(u)}{\pi_m(u)}\\
&=
-H(\nu)-\sum_{u\in X_m}\nu(u)\log\pi_m(u)\\
&=
m\log\varphi-H(\nu)-\log\pi(0)
{}+\bigl(\eta_{11}(\nu)-\eta_{00}(\nu)\bigr)\log\varphi.
\end{aligned}
\]
Since \(\eta_{11}(\nu)-\eta_{00}(\nu)\in[-1,1]\), the two-sided bound follows.
\end{proof}

\begin{remark}[Specificity of the endpoint collapse]\label{rem:endpoint-collapse-specificity}
Theorem~\ref{thm:parry-kl-identity} depends essentially on
Proposition~\ref{prop:parry-endpoint-collapse}. For a general shift of finite
type, Parry cylinder weights are controlled by boundary states in an edge
presentation but need not collapse to a scalar endpoint statistic of the kind
available here. The exact identity is part of what makes the golden-mean
setting especially rigid.
\end{remark}

\begin{corollary}[No fold-entropy loss on the Sturmian slice]\label{cor:no-fold-entropy-loss}
Assume \(\beta=\alpha\). Then
\[
\KL(\pi_m^{\mathrm{fold}}\|\pi_m)
=
m\log\varphi-H(\mu_m)-\log\pi(0)
{}+\bigl(\eta_{11}(\pi_m^{\mathrm{fold}})
-\eta_{00}(\pi_m^{\mathrm{fold}})\bigr)\log\varphi.
\]
Likewise, for every \(N\ge 1\),
\[
\KL(\hat\pi_{m,N}\|\pi_m)
=
m\log\varphi-H(\hat\mu_{m,N})-\log\pi(0)
{}+\bigl(\eta_{11}(\hat\pi_{m,N})
-\eta_{00}(\hat\pi_{m,N})\bigr)\log\varphi.
\]
\end{corollary}

\begin{proof}
Apply Theorem~\ref{thm:parry-kl-identity} with
\(\nu=\pi_m^{\mathrm{fold}}\) and \(\nu=\hat\pi_{m,N}\), and use
Corollary~\ref{cor:sturmian-exact-relabeling} to replace
\(H(\pi_m^{\mathrm{fold}})\) by \(H(\mu_m)\) and
\(H(\hat\pi_{m,N})\) by \(H(\hat\mu_{m,N})\).
\end{proof}

\begin{proposition}[Order-uniform R\'enyi window relative to the Parry measure]
\label{prop:parry-renyi-window}
For every probability law \(\nu\) on \(X_m\) and every \(q\in[0,\infty]\),
\[
m\log\varphi-H_q(\nu)-\log\pi(0)-\log\varphi
\le
D_q(\nu\|\pi_m)
\le
m\log\varphi-H_q(\nu)-\log\pi(0)+\log\varphi.
\]
\end{proposition}

\begin{proof}
For \(q=1\), this is Theorem~\ref{thm:parry-kl-identity}. For \(q=0\), let
\(S:=\supp(\nu)\). Proposition~\ref{prop:parry-endpoint-collapse} gives
\[
\pi(0)\varphi^{-(m+1)}
\le
\pi_m(u)
\le
\pi(0)\varphi^{-(m-1)}
\qquad (u\in X_m),
\]
and therefore
\[
|S|\pi(0)\varphi^{-(m+1)}
\le
\pi_m(S)
\le
|S|\pi(0)\varphi^{-(m-1)}.
\]
Since \(D_0(\nu\|\pi_m)=-\log \pi_m(S)\) and \(H_0(\nu)=\log|S|\), the
displayed bounds follow.

For \(q=\infty\), the same pointwise bounds on \(\pi_m(u)\) imply
\[
\frac{\varphi^{m-1}}{\pi(0)}\max_{u\in X_m}\nu(u)
\le
\max_{u:\,\nu(u)>0}\frac{\nu(u)}{\pi_m(u)}
\le
\frac{\varphi^{m+1}}{\pi(0)}\max_{u\in X_m}\nu(u).
\]
Taking logarithms and using \(H_\infty(\nu)=-\log\max_u \nu(u)\) gives the
claim.

Now fix \(q\in(0,\infty)\setminus\{1\}\). By
Proposition~\ref{prop:parry-endpoint-collapse},
\[
\pi_m(u)
=
\pi(0)\varphi^{-m+\theta(u)},
\qquad
\theta(u):=\mathbf{1}_{\{u_1=0\}}-\mathbf{1}_{\{u_m=1\}}\in[-1,1].
\]
Hence
\[
\sum_{u\in X_m}\nu(u)^q\pi_m(u)^{1-q}
=
\pi(0)^{1-q}\varphi^{m(q-1)}
\sum_{u\in X_m}\nu(u)^q\varphi^{(1-q)\theta(u)}.
\]
Because \(\theta(u)\in[-1,1]\), one has
\[
\varphi^{-|1-q|}\sum_{u\in X_m}\nu(u)^q
\le
\sum_{u\in X_m}\nu(u)^q\varphi^{(1-q)\theta(u)}
\le
\varphi^{|1-q|}\sum_{u\in X_m}\nu(u)^q.
\]
Taking logarithms, dividing by \(q-1\), and using the definition of \(H_q\)
gives the same two-sided bound with error at most \(\log\varphi\).
\end{proof}

\subsection{Support Conditioning and Endpoint Class Counts}

\begin{proposition}[Support conditioning for R\'enyi divergence]
\label{prop:renyi-support-conditioning}
Let \(E\) be finite, let \(\rho\) be a probability law on \(E\), and let
\(S\subset \supp(\rho)\) satisfy \(\rho(S)>0\). Writing
\(\rho^S:=\rho(\cdot\mid S)\), every probability law \(\nu\) on \(E\) with
\(\supp(\nu)\subset S\) satisfies
\[
D_q(\nu\|\rho)
=
D_q(\nu\|\rho^S)-\log \rho(S)
\qquad (q\in[0,\infty]).
\]
Consequently,
\[
\inf_{\substack{\nu\in\Delta(E)\\ \supp(\nu)\subset S}} D_q(\nu\|\rho)
=
-\log \rho(S)
\qquad (q\in[0,\infty]).
\]
If \(q\in(0,\infty]\), the minimizer is unique and equals \(\rho^S\). If
\(q=0\), the minimizers are exactly the laws whose support is \(S\).
\end{proposition}

\begin{proof}
For \(q=1\), the identity is the usual conditioning formula for relative
entropy. For \(q=0\),
\[
D_0(\nu\|\rho)
=
-\log \rho(\supp(\nu))
=
-\log \rho^S(\supp(\nu))-\log \rho(S)
=
D_0(\nu\|\rho^S)-\log \rho(S).
\]
For \(q=\infty\),
\[
D_\infty(\nu\|\rho)
=
\log\max_{x:\,\nu(x)>0}\frac{\nu(x)}{\rho(x)}
=
\log\max_{x:\,\nu(x)>0}\frac{\nu(x)}{\rho(S)\rho^S(x)}
=
D_\infty(\nu\|\rho^S)-\log \rho(S).
\]
Finally, for \(q\in(0,\infty)\setminus\{1\}\),
\[
\sum_{x\in E}\nu(x)^q\rho(x)^{1-q}
=
\sum_{x\in S}\nu(x)^q\bigl(\rho(S)\rho^S(x)\bigr)^{1-q}
=
\rho(S)^{1-q}\sum_{x\in S}\nu(x)^q(\rho^S(x))^{1-q},
\]
which gives the claimed identity after taking logarithms and dividing by
\(q-1\).

For every \(q>0\), the standard nonnegativity \(D_q(\nu\|\rho^S)\ge 0\) shows
that \(D_q(\nu\|\rho)\ge -\log\rho(S)\). Equality holds exactly when
\(D_q(\nu\|\rho^S)=0\), namely when \(\nu=\rho^S\). For \(q=0\),
\[
D_0(\nu\|\rho^S)=-\log \rho^S(\supp(\nu))\ge 0,
\]
with equality if and only if \(\rho^S(\supp(\nu))=1\), equivalently
\(\supp(\nu)=S\).
\end{proof}

\begin{lemma}[Endpoint class counts]\label{lem:endpoint-class-counts}
Extend the Fibonacci sequence by \(F_0:=0\), \(F_1:=1\), and
\(F_{n+2}=F_{n+1}+F_n\). For \(m\ge 2\), let
\[
X_m^{ij}:=\{u\in X_m:u_1=i,\ u_m=j\},
\qquad i,j\in\{0,1\}.
\]
Then
\[
|X_m^{00}|=F_m,\qquad
|X_m^{01}|=|X_m^{10}|=F_{m-1},\qquad
|X_m^{11}|=F_{m-2}.
\]
\end{lemma}

\begin{proof}
The admissible length-\(m\) words in \(X_m\) are the paths of length \(m-1\) in
the golden-mean graph with adjacency matrix
\[
A=\begin{pmatrix}1&1\\1&0\end{pmatrix}.
\]
For the standard Fibonacci extension above,
\[
A^n=
\begin{pmatrix}
F_{n+1} & F_n\\
F_n & F_{n-1}
\end{pmatrix}
\qquad (n\ge 1).
\]
Hence the number of admissible words of length \(m\) with endpoints \((i,j)\) is
\((A^{m-1})_{ij}\), which gives the displayed formulas.
\end{proof}

\subsection{Sharp Support Barriers}

\begin{theorem}[Sharp support barrier against the Parry measure]\label{thm:parry-support-barrier}
Let \(m\ge 2\) and \(1\le k\le F_{m+2}\). Define
\[
M_m(k):=\max\{\pi_m(S):S\subset X_m,\ |S|\le k\}.
\]
Then
\[
M_m(k)=
\begin{cases}
k\,\pi(0)\varphi^{-(m-1)},
& 1\le k\le F_m,\\[0.4em]
F_m\pi(0)\varphi^{-(m-1)}+(k-F_m)\pi(0)\varphi^{-m},
& F_m<k\le F_m+2F_{m-1},\\[0.4em]
1-(F_{m+2}-k)\pi(0)\varphi^{-(m+1)},
& F_m+2F_{m-1}<k\le F_{m+2}.
\end{cases}
\]
Moreover,
\[
\inf_{\substack{\nu\in\Delta(X_m)\\ |\supp(\nu)|\le k}}D_q(\nu\|\pi_m)
=
-\log M_m(k),
\qquad q\in[0,\infty],
\]
and
\[
\inf_{\substack{\nu\in\Delta(X_m)\\ |\supp(\nu)|\le k}}\TV(\nu,\pi_m)
=
1-M_m(k).
\]
If \(q\in(0,\infty]\), the R\'enyi infimum is attained uniquely by conditioning
\(\pi_m\) on a maximizing set \(S\subset X_m\); equivalently, one takes all
words from the heaviest endpoint classes until the budget \(k\) is exhausted.
If \(q=0\), the minimizers are exactly the laws whose support is such a
maximizing set. The same conditioned Parry laws attain the TV infimum.
\end{theorem}

\begin{proof}
The argument has three layers. One first orders individual Parry cylinder
weights by endpoint class, then maximizes the mass of a \(k\)-point set by a
greedy filling of these endpoint classes, and finally invokes
Proposition~\ref{prop:renyi-support-conditioning} to identify the optimal
R\'enyi law on a fixed support.

By Proposition~\ref{prop:parry-endpoint-collapse}, every word in \(X_m^{00}\) has
Parry weight \(\pi(0)\varphi^{-(m-1)}\), every word in \(X_m^{01}\cup X_m^{10}\)
has weight \(\pi(0)\varphi^{-m}\), and every word in \(X_m^{11}\) has weight
\(\pi(0)\varphi^{-(m+1)}\). These three values are strictly ordered. Together
with Lemma~\ref{lem:endpoint-class-counts}, this implies that the largest
possible Parry mass of a \(k\)-point subset is obtained greedily by exhausting
the endpoint classes in the order
\[
X_m^{00},\qquad X_m^{01}\cup X_m^{10},\qquad X_m^{11},
\]
which gives the formula for \(M_m(k)\).

Now let \(S\subset X_m\) and let \(\nu\) be supported on \(S\). By
Proposition~\ref{prop:renyi-support-conditioning},
\[
D_q(\nu\|\pi_m)
\ge
-\log \pi_m(S)
\qquad (q\in[0,\infty]),
\]
with equality for the conditioned law \(\pi_m^S:=\pi_m(\cdot\mid S)\), and
with uniqueness when \(q\in(0,\infty]\). Therefore
\[
\inf_{\substack{\nu\in\Delta(X_m)\\ |\supp(\nu)|\le k}}D_q(\nu\|\pi_m)
=
-\log\max_{\substack{S\subset X_m\\ |S|\le k}}\pi_m(S)
=
-\log M_m(k)
\]
for every \(q\in[0,\infty]\). The characterization of minimizers follows from
the same proposition.

For total variation, if \(\nu\) is supported on \(S\), then
\[
\TV(\nu,\pi_m)
\ge
\pi_m(S^c)
=
1-\pi_m(S).
\]
Equality holds for \(\nu=\pi_m^S\), since then
\[
\TV(\pi_m^S,\pi_m)
=
\frac12\sum_{u\in X_m}\bigl|\pi_m^S(u)-\pi_m(u)\bigr|
=
1-\pi_m(S).
\]
Taking the best support set \(S\) gives the total-variation formula.
\end{proof}

\begin{corollary}[Exact complexity-divergence law]\label{cor:exact-complexity-divergence-law}
Let \(m\ge 2\) and \(1\le k\le F_m\). Then
\[
\inf_{\substack{\nu\in\Delta(X_m)\\ |\supp(\nu)|\le k}}D_q(\nu\|\pi_m)
=
(m-1)\log\varphi-\log\pi(0)-\log k,
\qquad q\in[0,\infty],
\]
and
\[
\inf_{\substack{\nu\in\Delta(X_m)\\ |\supp(\nu)|\le k}}\TV(\nu,\pi_m)
=
1-k\pi(0)\varphi^{-(m-1)}.
\]
If \(q\in(0,\infty]\), the R\'enyi minimizers are exactly the uniform laws on
\(k\)-point subsets of \(X_m^{00}\). If \(q=0\), the minimizers are exactly the
laws whose support is such a \(k\)-point subset. The TV minimizers are the
conditioned Parry laws on these same subsets.

Consequently, if \(k_m\le F_m\) and
\[
\lim_{m\to\infty}\frac{1}{m}\log k_m=\gamma\in[0,\log\varphi],
\]
then
\[
\lim_{m\to\infty}
\inf_{\substack{\nu\in\Delta(X_m)\\ |\supp(\nu)|\le k_m}}
\frac{1}{m}D_q(\nu\|\pi_m)
=
\log\varphi-\gamma
=
h_{\mathrm{top}}(X)-\gamma,
\qquad q\in[0,\infty].
\]
In particular, every subexponential support budget still incurs the full rate
\(\log\varphi\), and cancelling the Parry penalty requires support growth of
order \(\varphi^m\).
\end{corollary}

\begin{proof}
The first branch of Theorem~\ref{thm:parry-support-barrier} gives the displayed
formulas. The description of the minimizers follows because all words in
\(X_m^{00}\) carry the same Parry weight \(\pi(0)\varphi^{-(m-1)}\), so the
conditioned Parry law on any \(k\)-point subset of \(X_m^{00}\) is uniform.
The limit statement is obtained by dividing the R\'enyi formula by \(m\) and
letting \(m\to\infty\).
\end{proof}

\section{Injective Placement into Primitive Shifts of Finite Type}
\label{sec:injective-placement}

\subsection{Optimal Injective Relabeling of the Realized Sturmian Law}

\begin{theorem}[Optimal injective relabeling of the realized Sturmian law]
\label{thm:optimal-injective-relabeling}
Assume \(\beta=\alpha\) and \(m\ge 6\). Let \(S_m(\alpha)\) be the realized
Sturmian language of length \(m\), let \(\mu_m\) be its factor law, and define
\[
\mathcal J_m
:=
\{\iota:S_m(\alpha)\hookrightarrow X_m:\ \iota\text{ injective}\},
\qquad
L_q(m)
:=
\inf_{\iota\in\mathcal J_m}D_q(\iota_*\mu_m\|\pi_m).
\]
Then for every \(q\in[0,\infty]\),
\[
L_q(m)
=
(m-1)\log\varphi-\log\pi(0)-H_q(\mu_m).
\]
Every injection whose image is contained in \(X_m^{00}\) attains the infimum.
Moreover, the canonical folded law satisfies
\[
0\le D_q(\pi_m^{\mathrm{fold}}\|\pi_m)-L_q(m)\le 2\log\varphi,
\qquad q\in[0,\infty],
\]
and for \(q=1\),
\[
\KL(\pi_m^{\mathrm{fold}}\|\pi_m)-L_1(m)
=
\bigl(1+\eta_{11}(\pi_m^{\mathrm{fold}})
-\eta_{00}(\pi_m^{\mathrm{fold}})\bigr)\log\varphi.
\]
\end{theorem}

\begin{proof}
Let
\[
c_m:=\pi(0)\varphi^{-(m-1)}.
\]
By Proposition~\ref{prop:parry-endpoint-collapse}, every word in \(X_m^{00}\)
has Parry mass \(c_m\), and no word in \(X_m\) has larger Parry mass. By
Lemma~\ref{lem:endpoint-class-counts},
\[
|X_m^{00}|=F_m.
\]
Since \(|S_m(\alpha)|=m+1\) and \(F_m\ge m+1\) for \(m\ge 6\), there exist
injections \(\iota\in\mathcal J_m\) with \(\iota(S_m(\alpha))\subset X_m^{00}\).

Fix such an injection and write \(S:=\iota(S_m(\alpha))\). Then \(|S|=m+1\),
\(\pi_m(S)=(m+1)c_m\), and the conditioned Parry law \(\pi_m^S\) is uniform on
\(S\). Because \(\iota_*\mu_m\) is a relabeling of \(\mu_m\), it has the same
R\'enyi entropies. Therefore
\[
D_q(\iota_*\mu_m\|\pi_m^S)=\log|S|-H_q(\mu_m)
\qquad (q\in[0,\infty]),
\]
and Proposition~\ref{prop:renyi-support-conditioning} gives
\[
D_q(\iota_*\mu_m\|\pi_m)
=
\log(m+1)-H_q(\mu_m)-\log((m+1)c_m)
=
-\log c_m-H_q(\mu_m).
\]
This proves the upper bound
\[
L_q(m)\le -\log c_m-H_q(\mu_m)
=(m-1)\log\varphi-\log\pi(0)-H_q(\mu_m).
\]

For the matching lower bound, fix an arbitrary \(\iota\in\mathcal J_m\), write
\(p_u:=\mu_m(u)\), and set \(\rho_u:=\pi_m(\iota(u))\). Since \(\rho_u\le c_m\)
for every \(u\in S_m(\alpha)\), the Shannon case satisfies
\[
D_1(\iota_*\mu_m\|\pi_m)
=
\sum_{u\in S_m(\alpha)}p_u\log\frac{p_u}{\rho_u}
\ge
\sum_{u\in S_m(\alpha)}p_u\log\frac{p_u}{c_m}
=
-H_1(\mu_m)-\log c_m.
\]
For \(q\in(0,\infty)\setminus\{1\}\), monotonicity of \(x\mapsto x^{1-q}\)
implies
\[
\frac{1}{q-1}\log\!\Bigl(\sum_{u\in S_m(\alpha)}p_u^q\rho_u^{1-q}\Bigr)
\ge
\frac{1}{q-1}\log\!\Bigl(c_m^{1-q}\sum_{u\in S_m(\alpha)}p_u^q\Bigr),
\]
because \(q>1\) gives a direct lower bound inside the logarithm while
\(q<1\) gives an upper bound inside the logarithm and a negative prefactor.
Hence
\[
D_q(\iota_*\mu_m\|\pi_m)\ge -\log c_m-H_q(\mu_m)
\qquad \bigl(q\in(0,\infty)\setminus\{1\}\bigr).
\]
For \(q=\infty\),
\[
D_\infty(\iota_*\mu_m\|\pi_m)
=
\log\max_{u\in S_m(\alpha)}\frac{p_u}{\rho_u}
\ge
\log\max_{u\in S_m(\alpha)}\frac{p_u}{c_m}
=
-\log c_m-H_\infty(\mu_m).
\]
Finally, since \(|S_m(\alpha)|=m+1\),
\[
\pi_m(\iota(S_m(\alpha)))\le (m+1)c_m=c_m e^{H_0(\mu_m)},
\]
and therefore
\[
D_0(\iota_*\mu_m\|\pi_m)
=
-\log \pi_m(\iota(S_m(\alpha)))
\ge
-\log c_m-H_0(\mu_m).
\]
The lower bound is now proved for every \(q\in[0,\infty]\), so the formula for
\(L_q(m)\) follows. Every injection into \(X_m^{00}\) attains equality by the
first part of the argument.

The left inequality
\[
0\le D_q(\pi_m^{\mathrm{fold}}\|\pi_m)-L_q(m)
\]
is immediate because \(\pi_m^{\mathrm{fold}}=(B_m)_*\mu_m\) with
\(B_m\in\mathcal J_m\) by Corollary~\ref{cor:sturmian-exact-relabeling}. The
upper bound follows from the R\'enyi invariance
\[
H_q(\pi_m^{\mathrm{fold}})=H_q(\mu_m)
\]
and Proposition~\ref{prop:parry-renyi-window}, which give
\[
D_q(\pi_m^{\mathrm{fold}}\|\pi_m)
\le
m\log\varphi-H_q(\mu_m)-\log\pi(0)+\log\varphi
=
L_q(m)+2\log\varphi.
\]
For \(q=1\), Corollary~\ref{cor:no-fold-entropy-loss} yields
\[
\KL(\pi_m^{\mathrm{fold}}\|\pi_m)
=
m\log\varphi-H_1(\mu_m)-\log\pi(0)
+\bigl(\eta_{11}(\pi_m^{\mathrm{fold}})
-\eta_{00}(\pi_m^{\mathrm{fold}})\bigr)\log\varphi,
\]
so subtracting the displayed formula for \(L_1(m)\) gives the exact identity.
\end{proof}

\begin{remark}[Injective relabeling versus support budgets]
\label{rem:injective-relabeling-stronger}
Theorem~\ref{thm:optimal-injective-relabeling} does more than constrain the
image support size. It fixes the realized Sturmian weight vector \(\mu_m\) and
allows only injective placements of those \(m+1\) probabilities inside the
ambient language \(X_m\). The resulting optimum is therefore a coding
invariance statement for the realized source, not just a best \((m+1)\)-state
approximation problem over arbitrary laws on \(X_m\). The finite exceptional
cases \(m<6\) are irrelevant for the asymptotic corollaries below and are
absorbed there into the \(O_\alpha(1)\) constants.
\end{remark}

\subsection{Exact Escort-Pressure Defect Formula}

The bounded additive defect in
Theorem~\ref{thm:optimal-injective-relabeling} admits a completely explicit
description for every R\'enyi order. The defect is not an entropy residual
but a boundary cocycle evaluated against the escort distribution of the
folded law.

\begin{theorem}[Full R\'enyi defect formula for the canonical folded law]
\label{thm:full-renyi-defect-formula}
Let
\[
\eta_m:=\pi_m^{\mathrm{fold}},\qquad Y_m:=\supp(\eta_m)\subset X_m.
\]
For \(x\in Y_m\), define the endpoint statistic
\[
\vartheta(x):=\mathbf{1}_{\{x_1=0\}}-\mathbf{1}_{\{x_m=1\}}\in\{-1,0,1\},
\]
the boundary cost
\[
c(x):=(1-\vartheta(x))\log\varphi\in\{0,\log\varphi,2\log\varphi\},
\]
and write
\[
Y_m^{ij}:=\{x\in Y_m:x_1=i,\ x_m=j\},\qquad i,j\in\{0,1\}.
\]
For \(q\in(0,\infty)\), define the \(q\)-escort law of \(\eta_m\) by
\[
\mathsf{E}_{m,q}(x):=\frac{\eta_m(x)^q}{\sum_{y\in Y_m}\eta_m(y)^q},
\qquad x\in Y_m.
\]
Then for each \(q\in(0,\infty)\setminus\{1\}\),
\begin{equation}\label{eq:escort-defect}
D_q(\eta_m\|\pi_m)-L_q(m)
=
\frac{1}{q-1}
\log\!\Bigl(
\mathsf{E}_{m,q}(Y_m^{00})
+\varphi^{\,q-1}\,\mathsf{E}_{m,q}(Y_m^{01}\cup Y_m^{10})
+\varphi^{\,2(q-1)}\,\mathsf{E}_{m,q}(Y_m^{11})
\Bigr).
\end{equation}
The \(q=1\) limit recovers the exact constant of
Theorem~\ref{thm:optimal-injective-relabeling}:
\begin{equation}\label{eq:defect-q1}
D_1(\eta_m\|\pi_m)-L_1(m)
=
\bigl(1+\eta_{11}(\eta_m)-\eta_{00}(\eta_m)\bigr)\log\varphi.
\end{equation}
The \(q=0\) endpoint formula is
\begin{equation}\label{eq:defect-q0}
D_0(\eta_m\|\pi_m)-L_0(m)
=
-\log\frac{|Y_m^{00}|+\varphi^{-1}(|Y_m^{01}|+|Y_m^{10}|)
+\varphi^{-2}|Y_m^{11}|}{|Y_m|}.
\end{equation}
The \(q=\infty\) endpoint formula is
\begin{equation}\label{eq:defect-qinfty}
D_\infty(\eta_m\|\pi_m)-L_\infty(m)
=
\log\max_{x\in Y_m}
\Bigl(
\varphi^{\,1-\vartheta(x)}
\frac{\eta_m(x)}{\max_{y\in Y_m}\eta_m(y)}
\Bigr).
\end{equation}
In particular, for all \(q\in[0,\infty]\),
\[
0\le D_q(\eta_m\|\pi_m)-L_q(m)\le 2\log\varphi,
\]
and the defect is entirely determined by the distribution of the folded image
across the endpoint classes under the \(q\)-escort geometry.
\end{theorem}

\begin{proof}
By Proposition~\ref{prop:parry-endpoint-collapse}, for any \(x\in X_m\),
\[
\pi_m(x)=\pi(0)\varphi^{-m+\vartheta(x)}.
\]
For \(q\in(0,\infty)\setminus\{1\}\), write
\(Z_{m,q}:=\sum_{x\in Y_m}\eta_m(x)^q\). Then
\[
\begin{aligned}
D_q(\eta_m\|\pi_m)
&=
\frac{1}{q-1}
\log\sum_{x\in Y_m}\eta_m(x)^q
\bigl(\pi(0)\varphi^{-m+\vartheta(x)}\bigr)^{1-q}\\
&=
m\log\varphi-\log\pi(0)
+\frac{1}{q-1}
\log\sum_{x\in Y_m}\eta_m(x)^q\varphi^{(1-q)\vartheta(x)}\\
&=
m\log\varphi-\log\pi(0)-H_q(\eta_m)
+\frac{1}{q-1}
\log\sum_{x\in Y_m}\mathsf{E}_{m,q}(x)\,\varphi^{(1-q)\vartheta(x)}.
\end{aligned}
\]
Since \(H_q(\eta_m)=H_q(\mu_m)\) by
Corollary~\ref{cor:sturmian-exact-relabeling} and
\(L_q(m)=(m-1)\log\varphi-\log\pi(0)-H_q(\mu_m)\), subtraction gives
\[
D_q(\eta_m\|\pi_m)-L_q(m)
=
\log\varphi
+\frac{1}{q-1}
\log\sum_{x\in Y_m}\mathsf{E}_{m,q}(x)\,\varphi^{(1-q)\vartheta(x)}.
\]
Absorbing \(\log\varphi\) into the logarithm yields
\[
D_q(\eta_m\|\pi_m)-L_q(m)
=
\frac{1}{q-1}
\log\sum_{x\in Y_m}\mathsf{E}_{m,q}(x)\,\varphi^{(q-1)(1-\vartheta(x))}.
\]
Since \(1-\vartheta(x)\) takes the value \(0\) on \(Y_m^{00}\), the value
\(1\) on \(Y_m^{01}\cup Y_m^{10}\), and the value \(2\) on \(Y_m^{11}\),
this is exactly~\eqref{eq:escort-defect}.

For the \(q=1\) limit, as \(q\to 1\) the escort law
\(\mathsf{E}_{m,q}\to\eta_m\), and
\[
\frac{\log\mathbb{E}(e^{(q-1)X})}{q-1}\to\mathbb{E}(X)
\qquad(q\to 1)
\]
with \(X=c(x)\). Hence
\[
D_1(\eta_m\|\pi_m)-L_1(m)
=\sum_{x\in Y_m}\eta_m(x)\,c(x)
=(\eta_{01}+\eta_{10}+2\eta_{11})\log\varphi.
\]
Using \(\eta_{00}+\eta_{01}+\eta_{10}+\eta_{11}=1\) gives
\(\eta_{01}+\eta_{10}+2\eta_{11}=1+\eta_{11}-\eta_{00}\),
recovering~\eqref{eq:defect-q1}.

For \(q=0\), one has \(D_0(\eta_m\|\pi_m)=-\log\pi_m(Y_m)\) and
\[
\pi_m(Y_m)
=
\pi(0)\varphi^{-(m-1)}|Y_m^{00}|
+\pi(0)\varphi^{-m}(|Y_m^{01}|+|Y_m^{10}|)
+\pi(0)\varphi^{-(m+1)}|Y_m^{11}|.
\]
Since \(L_0(m)=(m-1)\log\varphi-\log\pi(0)-\log|Y_m|\), subtraction
gives~\eqref{eq:defect-q0}.

For \(q=\infty\), let \(\eta_m^\star:=\max_{y\in Y_m}\eta_m(y)\). Then
\[
D_\infty(\eta_m\|\pi_m)
=
m\log\varphi-\log\pi(0)
+\log\max_{x\in Y_m}\eta_m(x)\,\varphi^{-\vartheta(x)}
\]
and
\(L_\infty(m)=(m-1)\log\varphi-\log\pi(0)+\log\eta_m^\star\).
Subtraction gives~\eqref{eq:defect-qinfty}.

Finally, \(c(x)\in[0,2\log\varphi]\) for all \(x\in Y_m\), so the defect
lies in \([0,2\log\varphi]\) for every \(q\in[0,\infty]\).
\end{proof}

\begin{remark}[Escort-pressure interpretation]
\label{rem:escort-pressure-interpretation}
Theorem~\ref{thm:full-renyi-defect-formula} shows that the entire
\(O(1)\) defect between the canonical folded law and the injective optimum
is neither an entropy residual nor a fold degeneracy, but the cumulative
generating function of the boundary cocycle \(c\) evaluated at the
\(q\)-escort distribution. In thermodynamic language,
\(D_q(\eta_m\|\pi_m)-L_q(m)\) is the \((q-1)\)-pressure of the endpoint
cost function \(c\) under the escort measure \(\mathsf{E}_{m,q}\). This
explains both the constant defect at \(q=1\) and the uniform boundedness
across all R\'enyi orders: the boundary cost takes only three values
\(\{0,\log\varphi,2\log\varphi\}\), so the pressure is uniformly controlled
by the span of the cocycle regardless of the escort geometry.
\end{remark}

\subsection{Bounded-Type Asymptotics}

\begin{proposition}[Bounded-type Sturmian R\'enyi entropy]\label{prop:bounded-type-sturmian-entropy}
Assume \(\beta=\alpha\), and suppose that \(\alpha\) has bounded partial
quotients. Then there exist
constants \(c_\alpha,C_\alpha>0\) such that for
every \(m\ge 1\) and every \(u\in S_m(\alpha)\),
\[
\frac{c_\alpha}{m}\le \mu_m(u)\le \frac{C_\alpha}{m}.
\]
Consequently,
\[
H_q(\mu_m)=\log m+O_\alpha(1)
\qquad (q\in[0,\infty]),
\]
with the \(O_\alpha(1)\) term uniform in \(q\).
\end{proposition}

\begin{proof}
Let \(p_n/q_n\) be the convergents of \(\alpha\), and choose \(n\) so that
\(q_n\le m<q_{n+1}\). By
Theorem~\ref{thm:spg-sturmian-cylinder-measure-continued-fraction-sandwich},
every length-\(m\) Sturmian cylinder satisfies
\[
\frac{1}{q_{n+1}+q_n}<\mu_m(u)<\frac{2}{q_n}.
\]
If the partial quotients of \(\alpha\) are bounded by \(A\), then
\(q_{n+1}\le (A+1)q_n\), and since \(q_n\le m<q_{n+1}\), one has
\[
\frac{1}{(A+2)m}\le \mu_m(u)\le \frac{2(A+1)}{m}.
\]
This proves the uniform \(m^{-1}\) bounds.

Because \(|S_m(\alpha)|=m+1\), one has
\[
H_0(\mu_m)=\log(m+1),
\qquad
H_\infty(\mu_m)=-\log\max_{u\in S_m(\alpha)}\mu_m(u)\ge \log m-\log C_\alpha.
\]
Since R\'enyi entropy on a finite alphabet is nonincreasing in \(q\),
\[
\log m-\log C_\alpha
\le
H_q(\mu_m)
\le
\log(m+1)
\qquad (q\in[0,\infty]).
\]
This yields the claimed uniform asymptotic.
\end{proof}

\begin{proposition}[Explicit bounded-type R\'enyi window]\label{prop:explicit-bounded-type-kl-window}
Assume \(\beta=\alpha\), and suppose that the partial quotients of \(\alpha\)
are bounded by \(A\). Then for every \(m\ge 1\) and every \(q\in[0,\infty]\),
\[
\log m-\log(2(A+1))
\le
H_q(\mu_m)
\le
\log(m+1).
\]
Consequently,
\[
m\log\varphi-\log(m+1)-\log\pi(0)-\log\varphi
\le
D_q(\pi_m^{\mathrm{fold}}\|\pi_m)
\]
\[
\le
m\log\varphi-\log m+\log(2(A+1))-\log\pi(0)+\log\varphi.
\]
In particular,
\[
D_q(\pi_m^{\mathrm{fold}}\|\pi_m)
-
\bigl(m\log\varphi-\log m\bigr)
\]
remains bounded by explicit constants depending only on \(A\) and \(\pi(0)\),
uniformly in \(q\).
\end{proposition}

\begin{proof}
The proof of Proposition~\ref{prop:bounded-type-sturmian-entropy} gives
\[
\frac{1}{(A+2)m}\le \mu_m(u)\le \frac{2(A+1)}{m}
\qquad (u\in S_m(\alpha)).
\]
Because \(|S_m(\alpha)|=m+1\), one has
\[
H_0(\mu_m)=\log(m+1),
\]
while the pointwise upper mass bound yields
\[
H_\infty(\mu_m)
=
-\log\max_{u\in S_m(\alpha)}\mu_m(u)
\ge
\log m-\log(2(A+1)).
\]
By monotonicity in \(q\),
\[
\log m-\log(2(A+1))
\le
H_q(\mu_m)
\le
\log(m+1)
\qquad (q\in[0,\infty]).
\]
Corollary~\ref{cor:sturmian-exact-relabeling} gives
\[
H_q(\pi_m^{\mathrm{fold}})=H_q(\mu_m)
\qquad (q\in[0,\infty]),
\]
and Proposition~\ref{prop:parry-renyi-window} then yields
\[
D_q(\pi_m^{\mathrm{fold}}\|\pi_m)
=
m\log\varphi-H_q(\mu_m)-\log\pi(0)+O(1),
\]
where the \(O(1)\) term is bounded in absolute value by \(\log\varphi\),
uniformly in \(q\). Combining these estimates gives the displayed window.
\end{proof}

\begin{corollary}[Injective-encoding invariance of the Parry barrier]
\label{cor:injective-relabeling-bounded-type}
Assume \(\beta=\alpha\), and suppose that \(\alpha\) has bounded partial
quotients. Then
\[
\inf_{\iota\in\mathcal J_m}D_q(\iota_*\mu_m\|\pi_m)
=
m\log\varphi-\log m+O_\alpha(1),
\qquad q\in[0,\infty],
\]
where the \(O_\alpha(1)\) term may be chosen uniformly in \(q\). Equivalently,
every injective encoding of the realized Sturmian law into \(X_m\) has the same
linear Parry penalty \(\log\varphi\) and the same logarithmic second-order term
\(-\log m\) up to bounded constants. In particular,
\[
D_q(\pi_m^{\mathrm{fold}}\|\pi_m)
=
\inf_{\iota\in\mathcal J_m}D_q(\iota_*\mu_m\|\pi_m)+O_\alpha(1),
\qquad q\in[0,\infty],
\]
and the additive error is in fact bounded between \(0\) and \(2\log\varphi\).
\end{corollary}

\begin{proof}
Theorem~\ref{thm:optimal-injective-relabeling} gives
\[
\inf_{\iota\in\mathcal J_m}D_q(\iota_*\mu_m\|\pi_m)
=
(m-1)\log\varphi-\log\pi(0)-H_q(\mu_m).
\]
Now apply Proposition~\ref{prop:bounded-type-sturmian-entropy} to substitute
\(H_q(\mu_m)=\log m+O_\alpha(1)\), uniformly in \(q\). The final claim is the
last part of Theorem~\ref{thm:optimal-injective-relabeling}.
\end{proof}

\begin{theorem}[Universal lossless coarse-graining envelope]
\label{thm:lossless-coarse-graining-envelope}
Assume \(\beta=\alpha\) and \(m\ge 6\). Let \(S_m(\alpha)\) be the realized
Sturmian language and \(\mu_m\) its factor law. Define the class of
deterministic maps that are injective on the realized source:
\[
\mathcal G_m
:=
\{G_m:\Omega_m\to X_m \text{ deterministic}:\
G_m|_{S_m(\alpha)}\text{ injective}\}.
\]
For \(G_m\in\mathcal G_m\), write
\(\nu_m^G:=(G_m)_*\mu_m\). Then for every \(q\in[0,\infty]\),
\[
\inf_{G_m\in\mathcal G_m}D_q(\nu_m^G\|\pi_m)
=
(m-1)\log\varphi-\log\pi(0)-H_q(\mu_m),
\]
and a map \(G_m\) attains the infimum if and only if
\(G_m(S_m(\alpha))\subset X_m^{00}\). Moreover, the canonical fold satisfies
\[
0\le D_q(\pi_m^{\mathrm{fold}}\|\pi_m)
-\inf_{G_m\in\mathcal G_m}D_q(\nu_m^G\|\pi_m)
\le 2\log\varphi.
\]
\end{theorem}

\begin{proof}
For any \(G_m\in\mathcal G_m\), the restriction
\(\iota:=G_m|_{S_m(\alpha)}:S_m(\alpha)\hookrightarrow X_m\) is injective, and
\(\nu_m^G=\iota_*\mu_m\). Thus the optimization over \(\mathcal G_m\) reduces
to the injective relabeling problem of
Theorem~\ref{thm:optimal-injective-relabeling}. The infimum, the
characterization of optimizers, and the bounded gap for the canonical fold all
follow from that theorem.
\end{proof}

\subsection{Primitive SFT Extension}

\begin{theorem}[Primitive SFT asymptotic R\'enyi barrier]
\label{thm:primitive-sft-renyi-barrier}
Let \(\Sigma_A\) be a primitive one-step shift of finite type with topological
entropy \(h=\log\lambda\). Write \(X_m^A\) for its admissible words of length
\(m\), and let \(\pi_m^A\) be the \(m\)-block law of its Parry measure. Then
there exist constants \(c_A,C_A>0\) such that
\[
c_A e^{-mh}\le \pi_m^A(w)\le C_A e^{-mh}
\qquad (w\in X_m^A).
\]
Consequently, if \((\nu_m)_{m\ge 1}\) is any family with
\(\nu_m\in\Delta(X_m^A)\) and
\[
\limsup_{m\to\infty}\frac{1}{m}\log |\supp(\nu_m)|\le \gamma<h,
\]
then
\[
\liminf_{m\to\infty}\frac{1}{m}D_q(\nu_m\|\pi_m^A)\ge h-\gamma
\qquad (q\in[0,\infty]).
\]
Conversely, for every \(\gamma\in[0,h]\), there exists a family
\((\nu_m)_{m\ge 1}\) with \(\nu_m\in\Delta(X_m^A)\) such that
\[
\frac{1}{m}\log |\supp(\nu_m)|\to \gamma,
\qquad
\frac{1}{m}D_q(\nu_m\|\pi_m^A)\to h-\gamma
\qquad (q\in[0,\infty]).
\]
\end{theorem}

\begin{proof}
Let \(u\) and \(v\) be positive left and right Perron-Frobenius eigenvectors
for \(A\), normalized so that \(\sum_i u_i v_i=1\). The standard Parry formula
for primitive one-step shifts \cite{LindMarcus1995,ParryPollicott1990Zeta}
gives
\[
\pi_m^A(i_0\cdots i_{m-1})=u_{i_0}v_{i_{m-1}}\lambda^{-(m-1)}
\]
for every admissible word \(i_0\cdots i_{m-1}\in X_m^A\). Since the alphabet is
finite and \(u_i v_j>0\) on admissible endpoint pairs, there are constants
\(c_A,C_A>0\) such that
\[
c_A e^{-mh}\le \pi_m^A(w)\le C_A e^{-mh}
\qquad (w\in X_m^A).
\]

Now let \(S_m:=\supp(\nu_m)\). By
Proposition~\ref{prop:renyi-support-conditioning},
\[
D_q(\nu_m\|\pi_m^A)\ge -\log \pi_m^A(S_m)
\qquad (q\in[0,\infty]).
\]
Using the upper block-mass bound,
\[
\pi_m^A(S_m)\le C_A |S_m| e^{-mh},
\]
so
\[
\frac{1}{m}D_q(\nu_m\|\pi_m^A)
\ge
h-\frac{1}{m}\log |S_m|-\frac{\log C_A}{m}.
\]
Taking \(\liminf\) proves the lower bound.

For sharpness, first fix \(\gamma\in[0,h)\) and set
\(k_m:=\max\{1,\lfloor e^{\gamma m}\rfloor\}\). Choose any subset
\(S_m\subset X_m^A\) with \(|S_m|=k_m\), and define
\(\nu_m:=\pi_m^A(\cdot\mid S_m)\). Then
\[
c_A k_m e^{-mh}
\le
\pi_m^A(S_m)
\le
C_A k_m e^{-mh},
\]
and Proposition~\ref{prop:renyi-support-conditioning} yields
\[
D_q(\nu_m\|\pi_m^A)=-\log \pi_m^A(S_m)=mh-\log k_m+O_A(1)
\qquad (q\in[0,\infty]).
\]
Since \(m^{-1}\log k_m\to \gamma\), this gives the desired limit. For
\(\gamma=h\), take \(\nu_m:=\pi_m^A\); then every \(D_q(\nu_m\|\pi_m^A)\) is
zero, while the language growth of a primitive shift satisfies
\(|X_m^A|=e^{hm+o(m)}\) \cite{LindMarcus1995}, so
\(m^{-1}\log |\supp(\nu_m)|\to h\).
\end{proof}

The asymptotic barrier of Theorem~\ref{thm:primitive-sft-renyi-barrier} is in
fact the coarsest shadow of a finite-\(m\) exact result. For every primitive
one-step SFT, the support-constrained R\'enyi optimization is already solved
exactly at the endpoint-class level. The golden-mean case is exceptional only
because its endpoint spectrum collapses to three Fibonacci classes, so the
general endpoint envelope reduces to a closed three-piece formula.

\begin{theorem}[Finite-\(m\) endpoint-class support envelope for primitive SFT]
\label{thm:primitive-sft-endpoint-envelope}
Adopt the notation of Theorem~\ref{thm:primitive-sft-renyi-barrier}.
Let \(\Gamma=\{\gamma_1>\cdots>\gamma_r\}\) be the distinct values taken by
\(u_iv_j\) over all admissible endpoint pairs \((i,j)\). For each
\(\ell\in\{1,\dots,r\}\), define the \(\ell\)-th endpoint class and its
capacity:
\[
B_{\ell,m}
:=
\{w\in X_m^A:u_{w_1}v_{w_m}=\gamma_\ell\},
\qquad
N_{\ell,m}:=|B_{\ell,m}|.
\]
Every \(w\in B_{\ell,m}\) has the same Parry mass
\(c_{\ell,m}:=\gamma_\ell\lambda^{-(m-1)}\). Write the cumulative capacities
\[
K_{j,m}:=\sum_{\ell=1}^j N_{\ell,m},
\qquad K_{0,m}:=0.
\]
For \(1\le k\le |X_m^A|\), let \(j=j(m,k)\) be the unique index satisfying
\(K_{j-1,m}<k\le K_{j,m}\), and define
\[
M_m^A(k)
:=
\sum_{\ell<j}N_{\ell,m}\,c_{\ell,m}
+
\bigl(k-K_{j-1,m}\bigr)\,c_{j,m}.
\]
Then for every \(q\in[0,\infty]\) and \(1\le k\le |X_m^A|\),
\[
\inf_{\substack{\nu\in\Delta(X_m^A)\\ |\supp(\nu)|\le k}}
D_q(\nu\|\pi_m^A)
=
-\log M_m^A(k),
\]
and
\[
\inf_{\substack{\nu\in\Delta(X_m^A)\\ |\supp(\nu)|\le k}}
\TV(\nu,\pi_m^A)
=
1-M_m^A(k).
\]
Any maximizing \(k\)-point support \(S\) must exhaust the heavier classes
\(B_{1,m},\dots,B_{j-1,m}\) before entering the frontier class \(B_{j,m}\).
For \(q\in(0,\infty]\), the R\'enyi minimizer is uniquely
\(\pi_m^A(\cdot\mid S)\); for \(q=0\), the minimizers are exactly the laws
whose support is a maximizing \(k\)-set.
\end{theorem}

\begin{proof}
Since \(c_{1,m}>\cdots>c_{r,m}>0\), the Parry mass of any \(k\)-point set
\(S\subset X_m^A\) is
\(\pi_m^A(S)=\sum_{\ell=1}^r|S\cap B_{\ell,m}|\,c_{\ell,m}\).
If \(S\) is mass-maximizing and some heavier class is not exhausted while a
lighter class is occupied, swapping one element from the lighter to the
heavier class strictly increases the mass. Hence every maximizing \(k\)-set
fills endpoint classes greedily in decreasing order of
\(c_{\ell,m}\), yielding \(\max_{|S|\le k}\pi_m^A(S)=M_m^A(k)\).
Proposition~\ref{prop:renyi-support-conditioning} then gives
\[
\inf_{|\supp(\nu)|\le k}D_q(\nu\|\pi_m^A)
=-\log\max_{|S|\le k}\pi_m^A(S)
=-\log M_m^A(k)
\]
for every \(q\in[0,\infty]\), and the TV formula follows from the same
argument as in Theorem~\ref{thm:parry-support-barrier}.
\end{proof}

\begin{remark}[Relation to the golden-mean case]
\label{rem:golden-mean-as-special-case}
For the golden-mean shift, the endpoint spectrum has exactly three levels
\(\gamma_1=\pi(0)\varphi,\;\gamma_2=\pi(0),\;\gamma_3=\pi(0)\varphi^{-1}\),
and the class capacities are \(N_{1,m}=F_m\),
\(N_{2,m}=2F_{m-1}\), \(N_{3,m}=F_{m-2}\) by
Lemma~\ref{lem:endpoint-class-counts}. The general envelope \(M_m^A(k)\)
of Theorem~\ref{thm:primitive-sft-endpoint-envelope} then reduces to the
three-piece formula of Theorem~\ref{thm:parry-support-barrier}. Thus the
golden-mean case is not special because it admits finite-\(m\) exact
optimization, but because its endpoint spectrum collapses to an explicitly
enumerable three-level structure with Fibonacci capacities.
\end{remark}

\begin{theorem}[Primitive SFT injective placement barrier]
\label{thm:primitive-sft-injective-barrier}
Let \(\Sigma_A\) be a primitive one-step shift of finite type with topological
entropy \(h=\log\lambda\). Let \(u,v\) be positive left and right
Perron-Frobenius eigenvectors for \(A\), normalized so that
\(\sum_i u_i v_i=1\). For an admissible word \(w=i_0\cdots i_{m-1}\), the
Parry formula gives
\(\pi_m^A(w)=u_{i_0}v_{i_{m-1}}\lambda^{-(m-1)}\). Define
\[
c_*:=\max_{i,j}u_i v_j,
\]
and let \((i_*,j_*)\) attain this maximum. Write
\[
\mathcal C_m(i_*,j_*):=\{w\in X_m^A: w_1=i_*,\;w_m=j_*\}.
\]
Then there exist \(\kappa_A>0\) and \(m_0\) such that for all \(m\ge m_0\),
\[
|\mathcal C_m(i_*,j_*)|\ge \kappa_A e^{mh}.
\]
Moreover, if \(\mu_m\) is any probability vector on a finite set \(S_m\), then
for every injection \(\iota:S_m\hookrightarrow X_m^A\) and every
\(q\in[0,\infty]\),
\[
D_q(\iota_*\mu_m\|\pi_m^A)
\ge
(m-1)h-\log c_*-H_q(\mu_m).
\]
If \(|S_m|\le |\mathcal C_m(i_*,j_*)|\), the lower bound is attained by any
injection into \(\mathcal C_m(i_*,j_*)\), and
\[
\inf_{\iota:S_m\hookrightarrow X_m^A}D_q(\iota_*\mu_m\|\pi_m^A)
=
(m-1)h-\log c_*-H_q(\mu_m).
\]
In particular, if \(\log|S_m|=o(m)\), then
\[
\inf_{\iota:S_m\hookrightarrow X_m^A}D_q(\iota_*\mu_m\|\pi_m^A)
=
mh-H_q(\mu_m)+O_A(1),
\]
where the \(O_A(1)\) term is uniform in \(q\).
\end{theorem}

\begin{proof}
By the Perron-Frobenius theorem,
\((A^{m-1})_{i_*j_*}=\lambda^{m-1}(v_{i_*}u_{j_*}+o(1))\), so there exist
\(\kappa_A>0\) and \(m_0\) with
\(|\mathcal C_m(i_*,j_*)|=(A^{m-1})_{i_*j_*}\ge \kappa_A\lambda^{m-1}\) for
\(m\ge m_0\).

Fix an injection \(\iota\) and write \(p_s=\mu_m(s)\),
\(\rho_s=\pi_m^A(\iota(s))\). By definition of \(c_*\),
\(\rho_s\le c_*\lambda^{-(m-1)}\) for all \(s\in S_m\).

For \(q=1\):
\(D_1=\sum_s p_s\log(p_s/\rho_s)
\ge \sum_s p_s\log(p_s/(c_*\lambda^{-(m-1)}))
=(m-1)h-\log c_*-H_1(\mu_m)\).

For \(q\in(0,\infty)\setminus\{1\}\): the bound
\(\rho_s\le c_*\lambda^{-(m-1)}\) gives
\(\rho_s^{1-q}\ge (c_*\lambda^{-(m-1)})^{1-q}\) when \(q>1\) and
\(\rho_s^{1-q}\le (c_*\lambda^{-(m-1)})^{1-q}\) when \(q<1\). In both cases
the prefactor \((q-1)^{-1}\) produces the same lower bound
\(D_q\ge (m-1)h-\log c_*-H_q(\mu_m)\).

For \(q=\infty\):
\(D_\infty=\log\max_s(p_s/\rho_s)\ge (m-1)h-\log c_*-H_\infty(\mu_m)\).

For \(q=0\):
\(\pi_m^A(\iota(S_m))\le |S_m|c_*\lambda^{-(m-1)}\), so
\(D_0=-\log\pi_m^A(\iota(S_m))\ge (m-1)h-\log c_*-H_0(\mu_m)\).

For attainability, if \(|S_m|\le |\mathcal C_m(i_*,j_*)|\), inject
\(S_m\) into \(\mathcal C_m(i_*,j_*)\). All image words then have identical
Parry mass \(c_*\lambda^{-(m-1)}\), so the lower bound is attained for
every \(q\). If \(\log|S_m|=o(m)\), then
\(|S_m|\le \kappa_A\lambda^{m-1}\) for large \(m\), and absorbing
\((m-1)h-\log c_*\) into \(mh+O_A(1)\) yields the last claim.
\end{proof}

\begin{corollary}[Low-complexity injective placement barrier]
\label{cor:low-complexity-source-barrier}
Let \(\Sigma_A\) be a primitive one-step SFT with topological entropy
\(h=\log\lambda\) and Parry \(m\)-block law \(\pi_m^A\). Let
\((S_m,\mu_m)_{m\ge 1}\) be a family of finite probability spaces satisfying
\[
|S_m|=m+O(1),
\qquad
H_q(\mu_m)=\log m+O(1)
\]
with the \(O(1)\) terms uniform in \(q\in[0,\infty]\). Then
\[
\inf_{\iota:S_m\hookrightarrow X_m^A}
D_q(\iota_{*}\mu_m\|\pi_m^A)
=
mh-\log m+O_A(1),
\qquad q\in[0,\infty],
\]
where \(O_A(1)\) is uniform in \(q\).
\end{corollary}

\begin{proof}
Since \(|S_m|=m+O(1)\), one has \(\log|S_m|=o(m)\).
Theorem~\ref{thm:primitive-sft-injective-barrier} therefore gives
\[
\inf_{\iota:S_m\hookrightarrow X_m^A}
D_q(\iota_{*}\mu_m\|\pi_m^A)
=
mh-H_q(\mu_m)+O_A(1),
\qquad q\in[0,\infty].
\]
Substituting \(H_q(\mu_m)=\log m+O(1)\), with the \(O(1)\) term uniform in
\(q\), yields
\[
\inf_{\iota:S_m\hookrightarrow X_m^A}
D_q(\iota_{*}\mu_m\|\pi_m^A)
=
mh-\log m+O_A(1),
\]
where the new \(O_A(1)\) remains uniform in \(q\).
\end{proof}

\section{Conjugacy Invariance and One Illustration}
\label{sec:conjugacy}

\begin{theorem}[Conjugacy invariance]
\label{thm:audit-conjugacy-invariance}
Let \(f:S^1\to S^1\) be an orientation-preserving circle homeomorphism
topologically conjugate to the irrational rotation
\(R_\alpha:x\mapsto x+\alpha\pmod 1\) via an orientation-preserving
homeomorphism \(h\) satisfying \(h\circ f=R_\alpha\circ h\). Fix an anchor
\(a\in S^1\) with \(h(a)=0\), and define the coding window
\(I_a:=[a,f(a))\). For \(x\in S^1\), let
\[
s_t^{(f)}(x):=\mathbf{1}_{I_a}(f^t(x)),
\qquad
s_t^{(\alpha)}(y):=\mathbf{1}_{[0,\alpha)}(y+t\alpha\bmod 1).
\]
Then for every \(x\in S^1\) and \(y:=h(x)\),
\[
s_t^{(f)}(x)=s_t^{(\alpha)}(y)
\qquad(t\ge 0).
\]
Consequently, for every \(m,N\ge 1\), the length-\(m\) microstate laws,
folded laws, empirical histograms, higher-block conjugacy, and all
R\'enyi and KL conclusions of Sections~\ref{sec:sturmian-rigidity}
through~\ref{sec:injective-placement} are identical for the two systems.
In particular, if \(f\) is a \(C^2\) orientation-preserving circle
diffeomorphism with irrational rotation number \(\alpha\), the Denjoy
theorem provides the required conjugacy, and the full theory applies.
\end{theorem}

\begin{proof}
Since \(h\) is orientation-preserving, the image of \(I_a=[a,f(a))\) is the
positively oriented arc from \(h(a)=0\) to
\(h(f(a))=R_\alpha(h(a))=\alpha\), that is, \(h(I_a)=[0,\alpha)\). For any
\(x\in S^1\) with \(y=h(x)\), the conjugacy gives
\(h(f^t(x))=R_\alpha^t(y)=y+t\alpha\pmod 1\), so
\[
s_t^{(f)}(x)
=\mathbf{1}_{I_a}(f^t(x))
=\mathbf{1}_{h(I_a)}(h(f^t(x)))
=\mathbf{1}_{[0,\alpha)}(y+t\alpha\bmod 1)
=s_t^{(\alpha)}(y).
\]
Since the symbol sequences agree pointwise for all \(t\), the length-\(m\)
sliding windows, and therefore all empirical counts, folded histograms, and
derived quantities, are identical.
\end{proof}

As a concrete illustration, consider the standard circle map
\[
f_{K,\Omega}(x)=x+\Omega-\frac{K}{2\pi}\sin(2\pi x)\pmod 1
\]
with \(K=0.3\) and \(\Omega\) tuned by bisection on the lift so that the
rotation number equals \(\varphi^{-1}\), giving
\(\Omega\approx 0.6166931270470741\). With anchor \(a=0.2\), coding window
\(I=[a,f(a))\), and initial point \(x_0=a\), the resulting symbolic sequence
is the canonical golden Sturmian coding in conjugacy coordinates.
At \(m=12\) and \(N=10^5\), the empirical folded histogram satisfies
\[
\TV(\hat\pi_{12,N},\pi_{12}^{\mathrm{fold}})
=2.43\times 10^{-5},
\qquad
\KL(\hat\pi_{12,N}\|\pi_{12}^{\mathrm{fold}})
=2.95\times 10^{-9},
\]
while the gaps to the Parry baseline remain
\[
\TV(\hat\pi_{12,N},\pi_{12})=0.9650,
\qquad
\KL(\hat\pi_{12,N}\|\pi_{12})=3.4443.
\]
The empirical law is already within \(10^{-5}\) total variation of its
deterministic limit, yet the structural divergence from the Parry measure
persists at order \(10^0\), consistent with
Corollary~\ref{cor:injective-relabeling-bounded-type}.

\section{Concluding Remarks}
\label{sec:concluding}

The canonical Zeckendorf fold provides a deterministic coarse-graining of
binary rotation words into the golden-mean shift language. The main
contribution of this paper is to show that the resulting divergence from the
Parry measure admits a complete finite-\(m\) analysis in which every quantity
is either closed-form or bounded by explicit constants.

The canonical Sturmian slice is the point at which the picture becomes
especially rigid. There the fold is not only injective on realized blocks; it
is conjugate to the classical higher-block presentation of the Sturmian source.
At the same time, the endpoint structure of the golden-mean Parry cylinders is
simple enough to make the Parry comparison exact.
Theorem~\ref{thm:optimal-injective-relabeling} shows that even the best
injective placement of the realized Sturmian law in \(X_m\) still has
divergence \(m\log\varphi-\log m+O_\alpha(1)\) for bounded-type slopes,
uniformly in the R\'enyi order. The folded law lies within a bounded additive
distance of this optimum.

This finite-\(m\) analysis calibrates a universal phenomenon: for any
primitive shift of finite type with topological entropy \(h\), any
deterministic coarse-graining injective on a subexponential-support source
incurs divergence at rate \(mh-\log m+O(1)\)
(Corollary~\ref{cor:low-complexity-source-barrier}). The golden-mean benchmark
is the unique setting in the Fibonacci family where the universal envelope
becomes exact, the optimizer is explicit, and the Fibonacci weighting is forced
by a contiguous numeration axiom
(Theorem~\ref{thm:fibonacci-weight-uniqueness}).

The limitations are also clear. The exact finite-\(m\) formulas rely on the
endpoint collapse geometry of the golden-mean shift. The injective-placement
problem is asymptotic in general, but the support-constrained Parry
optimization is finite-\(m\) exact for every primitive one-step SFT
(Theorem~\ref{thm:primitive-sft-endpoint-envelope}). What is special about the
golden-mean benchmark is therefore not exact solvability itself, but the
complete collapse of the endpoint spectrum to an explicitly enumerable
three-level structure. Whether analogous elementary closed forms exist for
other shifts of finite type with richer endpoint spectra, or for numeration
systems attached to other Pisot units, remains open. By topological conjugacy
(Theorem~\ref{thm:audit-conjugacy-invariance}), the entire theory transfers
verbatim to \(C^2\) circle diffeomorphisms with irrational rotation number.


\appendix
\section{Auxiliary Discrepancy and Empirical Entropy Bounds}
\label{app:discrepancy}

This appendix collects finite-sample control results that bound the empirical
folded histogram relative to its deterministic limit law. They complement the
structural Parry divergence analysis of the main text by showing that the
sampling error can be made arbitrarily small.

\subsection{Discrepancy of the Rotation Orbit}

The quantitative input is the one-dimensional star discrepancy of the
rotation orbit:
\[
D_N^*(\alpha;x_0)
:=
\sup_{0\le a\le 1}
\left|
\frac{1}{N}\sum_{t=0}^{N-1}\mathbf{1}_{[0,a)}(x_t)-a
\right|.
\]
This quantity is effective because length-$m$ rotation words are determined by a
finite partition of the circle into interval atoms.

\subsection{Total Variation Bounds}

\begin{theorem}[Discrepancy control of microstate and folded laws]\label{thm:tv-certificate-hist}
In the rotation model above, for every $m,N\ge 1$,
\[
\TV(\hat\mu_{m,N},\mu_m)\le 2m\,D_N^*(\alpha;x_0).
\]
Consequently,
\[
\TV(\hat\pi_{m,N},\pi_m^{\mathrm{fold}})
\le
2m\,D_N^*(\alpha;x_0).
\]
If \(\beta=\alpha\), then the sharper bounds
\[
\TV(\hat\mu_{m,N},\mu_m)\le (m+1)D_N^*(\alpha;x_0),
\qquad
\TV(\hat\pi_{m,N},\pi_m^{\mathrm{fold}})\le (m+1)D_N^*(\alpha;x_0)
\]
also hold.
\end{theorem}

\begin{proof}
For fixed $m$, the length-$m$ rotation words are constant on the interval atoms
of the partition of \(\TT\) cut out by the points
\[
0,\ \beta,\ -\alpha,\ \beta-\alpha,\ \dots,\ -(m-1)\alpha,\ \beta-(m-1)\alpha
\pmod 1.
\]
There are at most \(2m\) such endpoints and therefore at most \(2m\) interval
atoms. For any interval \(I=[a,b)\subset\TT\),
\[
\left|
\frac{1}{N}\sum_{t=0}^{N-1}\mathbf{1}_I(x_t)-\leb(I)
\right|
\le
2D_N^*(\alpha;x_0).
\]
Taking half the sum of the atomwise errors gives
\(\TV(\hat\mu_{m,N},\mu_m)\le 2mD_N^*(\alpha;x_0)\).

If \(\beta=\alpha\), the coding is Sturmian and the number of distinct
length-\(m\) factors is exactly \(m+1\)
\cite{MorseHedlund1940,LindMarcus1995}. Summing over the \(m+1\) nonempty
cylinder atoms yields the sharper constant. The folded estimates follow because
deterministic pushforward cannot increase total variation.
\end{proof}

\begin{corollary}[Optimal leading constant within the metallic-mean family]
\label{cor:metallic-mean-golden-unique-optimal-tv-constant}
For $A\ge 1$, let
\[
\alpha_A=[0;A,A,\dots],\qquad
\lambda_A:=\frac{A+\sqrt{A^2+4}}{2},\qquad
\kappa(A):=\frac{A}{\log\lambda_A}.
\]
Then there is a constant $C(A)<\infty$ such that for every $N\ge 2$ and every
$x_0\in\TT$,
\[
D_N^*(\alpha_A;x_0)\le \frac{\kappa(A)\log N+C(A)}{N}.
\]
Consequently, for every fixed resolution $m\ge 1$,
\[
\TV(\hat\pi_{m,N},\pi_m^{\mathrm{fold}})
\le
2m\frac{\kappa(A)\log N+C(A)}{N}.
\]
If \(\beta=\alpha_A\), this sharpens to
\[
\TV(\hat\pi_{m,N},\pi_m^{\mathrm{fold}})
\le
(m+1)\frac{\kappa(A)\log N+C(A)}{N}.
\]
Moreover, $\kappa(A)$ is strictly increasing on $[1,\infty)$, so the golden
branch $A=1$ uniquely minimizes the leading $(\log N)/N$ coefficient among
metallic means.
\end{corollary}

\begin{proof}
Let $p_n/q_n$ be the convergents of $\alpha_A$.  The standard discrepancy bound
\cite{KuipersNiederreiter1974} gives
\[
D_N^*(\alpha_A;x_0)\le \frac{3+nA}{N}
\qquad \text{whenever } q_n\le N<q_{n+1}.
\]
Since $q_{n+1}=Aq_n+q_{n-1}$ and $\lambda_A$ is the Perron root of
$x^2-Ax-1=0$, one has $q_n\ge \lambda_A^{n-1}$ for $n\ge 1$, hence
$n\le 1+\log N/\log\lambda_A$.  Substituting yields the discrepancy estimate
with $C(A)=3+A$.  The TV bound follows from
Theorem~\ref{thm:tv-certificate-hist}.  The monotonicity of
$\kappa(A)=A/\log\lambda_A$ on $[1,\infty)$ is elementary, so its minimum is
attained uniquely at $A=1$.
\end{proof}

\subsection{Cross-Resolution Residuals}

Let $r_{m+1\to m}:\Omega_{m+1}\to\Omega_m$ denote prefix truncation,
$r_{m+1\to m}(u_1\cdots u_{m+1})=u_1\cdots u_m$, and let
$\rho_{m+1\to m}:X_{m+1}\to X_m$ be the analogous truncation on folded words.
Define the projective residual
\[
E_{m,N}
:=
\TV\!\left(\hat\pi_{m,N},(\rho_{m+1\to m})_*\hat\pi_{m+1,N}\right).
\]

\begin{theorem}[Boundary law for the canonical fold]\label{prop:multiscale-residual-by-boundary}
For the canonical fold of Definition~\ref{def:canonical-fold}, let
\[
f_m(\omega):=\Fold_m(r_{m+1\to m}(\omega)),
\qquad
g_m(\omega):=\rho_{m+1\to m}(\Fold_{m+1}(\omega)),
\]
and define the boundary event set
\[
\mathcal{B}_{m+1}:=\{\omega\in\Omega_{m+1}: f_m(\omega)\neq g_m(\omega)\}.
\]
Then
\[
E_{m,N}\le \hat\mu_{m+1,N}(\mathcal{B}_{m+1}).
\]
Moreover, if \(\omega=u_1\cdots u_{m+1}\) and \(v:=u_1\cdots u_m\), then
\[
\omega\in\mathcal{B}_{m+1}
\iff
u_{m+1}=1\ \text{and}\ N(v)\ge F_{m+1}.
\]
\end{theorem}

\begin{proof}
Both $\hat\pi_{m,N}$ and $(\rho_{m+1\to m})_*\hat\pi_{m+1,N}$ are pushforwards
of $\hat\mu_{m+1,N}$, respectively by $f_m$ and $g_m$.  Therefore
\[
\TV((f_m)_*\nu,(g_m)_*\nu)\le \nu(f_m\neq g_m)
\]
for any probability measure $\nu$, and the first claim follows with
$\nu=\hat\mu_{m+1,N}$.

For the second claim, write \(N:=N(v)\). Since
\[
0\le N\le \sum_{k=1}^m F_{k+1}=F_{m+3}-2,
\]
the Zeckendorf expansion of \(N\) has the form
\[
Z(N)=(c_1,\dots,c_{m+1},0,0,\dots).
\]
If \(u_{m+1}=0\), then \(N(\omega)=N\), so \(f_m(\omega)=g_m(\omega)\).

Assume next that \(u_{m+1}=1\), so \(N(\omega)=N+F_{m+2}\). If \(N<F_{m+1}\),
then \(c_m=c_{m+1}=0\), hence
\[
Z(N+F_{m+2})=(c_1,\dots,c_m,1,0,0,\dots)
\]
is already a legal Zeckendorf expansion. Therefore
\[
g_m(\omega)=\rho_{m+1\to m}(\Fold_{m+1}(\omega))=(c_1,\dots,c_m)=f_m(\omega),
\]
so \(\omega\notin\mathcal{B}_{m+1}\).

Conversely, suppose \(u_{m+1}=1\) and \(\omega\notin\mathcal{B}_{m+1}\). Write
\[
Z(N+F_{m+2})=(d_1,\dots,d_{m+2},0,0,\dots).
\]
Since the first \(m\) digits agree, one has \(d_j=c_j\) for \(1\le j\le m\).
Subtracting the two Zeckendorf expansions gives
\[
F_{m+2}
=
(d_{m+1}-c_{m+1})F_{m+2}+d_{m+2}F_{m+3}.
\]
Because \(F_{m+3}>F_{m+2}\), necessarily \(d_{m+2}=0\) and
\(d_{m+1}-c_{m+1}=1\). Hence \(c_{m+1}=0\) and \(d_{m+1}=1\). Since
\((d_1,\dots,d_{m+2})\) is a legal Zeckendorf digit string and \(d_m=c_m\), one
must also have \(c_m=d_m=0\). Thus \(Z(N)\) has no nonzero digit in positions
\(m\) or \(m+1\), which is equivalent to \(N<F_{m+1}\).

We have proved
\[
\omega\notin\mathcal{B}_{m+1}
\iff
u_{m+1}=0\ \text{or}\ \bigl(u_{m+1}=1\ \text{and}\ N(v)<F_{m+1}\bigr),
\]
which is exactly the stated characterization.
\end{proof}

\begin{corollary}[Decomposition relative to the Parry measure]\label{cor:tv-decomposition-parry}
For every $m,N\ge 1$,
\[
\TV(\hat\pi_{m,N},\pi_m)
\le
2m\,D_N^*(\alpha;x_0)+\TV(\pi_m^{\mathrm{fold}},\pi_m).
\]
If \(\beta=\alpha\), one may replace \(2m\) by \(m+1\).
\end{corollary}

\begin{proof}
Theorem~\ref{thm:tv-certificate-hist} gives
$\TV(\hat\pi_{m,N},\pi_m^{\mathrm{fold}})\le 2mD_N^*(\alpha;x_0)$, and the
triangle inequality completes the proof.
\end{proof}

\subsection{Relative-Entropy Bounds}

\begin{lemma}[Relative entropy from total variation under strict positivity]\label{lem:kl-from-tv-qmin}
Let $p,q$ be probability measures on a finite set $S$, assume $p\ll q$, and
suppose that $q$ is strictly positive on $S$. Let
\[
q_{\min}:=\min_{s\in S}q(s)>0.
\]
Then
\[
\KL(p\|q)
\le
\log\!\left(1+\frac{\|p-q\|_1^2}{q_{\min}}\right)
\le
\frac{1}{q_{\min}}\|p-q\|_1^2
=
\frac{4}{q_{\min}}\TV(p,q)^2.
\]
\end{lemma}

\begin{proof}
The standard inequality \(\KL\le \log(1+\chi^2)\) gives
\[
\KL(p\|q)\le \log\bigl(1+\chi^2(p,q)\bigr),
\qquad
\chi^2(p,q):=\sum_{s\in S}\frac{(p(s)-q(s))^2}{q(s)}.
\]
Since \(q(s)\ge q_{\min}\) for every \(s\in S\),
\[
\chi^2(p,q)
\le
\frac{1}{q_{\min}}\sum_{s\in S}(p(s)-q(s))^2
\le
\frac{1}{q_{\min}}\|p-q\|_1^2.
\]
The second inequality uses \(\log(1+u)\le u\). See
\cite{SasonVerdu2016fDivergence}.
\end{proof}

\begin{definition}[Diophantine lower bound]
For a badly approximable irrational $\alpha$, let $c(\alpha)>0$ be a constant
such that
\[
\|k\alpha\|\ge \frac{c(\alpha)}{k}\qquad (k\ge 1),
\]
where $\|x\|:=\min_{n\in\mathbb{Z}}|x-n|$.
\end{definition}

\begin{proposition}[Lower bound on the smallest microstate atom]\label{prop:qmin-lowerbound-markov}
Assume $\beta=\alpha$ and let
\[
q_{\min}(m):=\min\{\mu_m(u):\mu_m(u)>0\}.
\]
Then
\[
q_{\min}(m)\ge \min_{1\le k\le m}\|k\alpha\|\ge \frac{c(\alpha)}{m}.
\]
For the golden slope $\alpha=\varphi^{-1}$ one may take
$c(\alpha)=1/\sqrt5$, hence
\[
q_{\min}(m)\ge \frac{1}{\sqrt5\,m}.
\]
\end{proposition}

\begin{proof}
When $\beta=\alpha$, the partition endpoints for length-$m$ microstates reduce
to the set $\{-k\alpha:0\le k\le m\}$.  Thus the positive atoms of $\mu_m$ are
exactly the adjacent gaps determined by these points, and the smallest one is
at least $\min_{1\le k\le m}\|k\alpha\|$.  The badly approximable lower bound
then yields $c(\alpha)/m$.  For $\varphi^{-1}$ the optimal Markov constant is
$1/\sqrt5$.
\end{proof}

\begin{corollary}[Microstate relative-entropy bound]\label{cor:rotation-microstate-kl-certificate}
Assume $\beta=\alpha$ and $\alpha$ is badly approximable.  Then
\[
\KL(\hat\mu_{m,N}\|\mu_m)
\le
\frac{4}{q_{\min}(m)}\TV(\hat\mu_{m,N},\mu_m)^2
\le
\frac{4(m+1)^2}{q_{\min}(m)}D_N^*(\alpha;x_0)^2.
\]
For $\alpha=\varphi^{-1}$ this becomes
\[
\KL(\hat\mu_{m,N}\|\mu_m)
\le
4\sqrt5\,m(m+1)^2 D_N^*(\alpha;x_0)^2.
\]
\end{corollary}

\begin{proof}
Let \(S_m:=\supp(\mu_m)\). In the Sturmian slice \(\beta=\alpha\), every
observed length-\(m\) word belongs to the Sturmian language, hence
\(\hat\mu_{m,N}\) and \(\mu_m\) may both be viewed as probability measures on
\(S_m\). On this set, \(\mu_m\) is strictly positive and
\[
\min_{u\in S_m}\mu_m(u)=q_{\min}(m).
\]
The claim therefore follows from Theorem~\ref{thm:tv-certificate-hist},
Lemma~\ref{lem:kl-from-tv-qmin}, and
Proposition~\ref{prop:qmin-lowerbound-markov}.
\end{proof}

\begin{lemma}[KL monotonicity under pushforward]\label{lem:kl-monotone-pushforward}
Let $F:S\to T$ be a map between finite sets and let $p,q$ be probability
measures on $S$ with $p\ll q$.  Then
\[
\KL(F_*p\|F_*q)\le \KL(p\|q).
\]
\end{lemma}

\begin{proof}
This is the data-processing inequality for relative entropy under deterministic
coarse graining.
\end{proof}

\begin{corollary}[Folded relative-entropy bound]\label{cor:rotation-folded-kl-certificate}
Under the assumptions of
Corollary~\ref{cor:rotation-microstate-kl-certificate},
\[
\KL(\hat\pi_{m,N}\|\pi_m^{\mathrm{fold}})
\le
\KL(\hat\mu_{m,N}\|\mu_m)
\le
\frac{4(m+1)^2}{q_{\min}(m)}D_N^*(\alpha;x_0)^2.
\]
\end{corollary}

\begin{proof}
Apply Lemma~\ref{lem:kl-monotone-pushforward} with $F=\Fold_m$.
\end{proof}

\section{Fiber-Uniform Lifts and I-Projections}\label{app:fiber-lift}

This appendix collects the microstate-level lift identities associated with the
canonical fold. They are used only to interpret the hidden degeneracy term in
Corollary~\ref{cor:sturmian-pure-degeneracy-ledger} and may be skipped on a
first reading.

\begin{definition}[Fiber multiplicity and uniform lift]\label{def:fiber-uniform-lift}
For \(x\in X_m\), define the fiber multiplicity
\[
d_m(x):=\bigl|\Fold_m^{-1}(x)\bigr|.
\]
For \(\eta\in\Delta(X_m)\), define its fiber-uniform lift
\[
(Q_m\eta)(a):=\frac{\eta(\Fold_m(a))}{d_m(\Fold_m(a))},
\qquad a\in\Omega_m.
\]
Then \((\Fold_m)_*(Q_m\eta)=\eta\). For \(\mu\in\Delta(\Omega_m)\), let
\[
\pi:=(\Fold_m)_*\mu\in\Delta(X_m),
\qquad
\mu^e:=Q_m\pi.
\]
\end{definition}

\begin{proposition}[Exact entropy ledger for the fold quotient]\label{prop:fold-entropy-ledger}
For every \(\mu\in\Delta(\Omega_m)\), with \(\pi=(\Fold_m)_*\mu\) and
\(\mu^e=Q_m\pi\),
\[
H(\mu)
=
H(\pi)+\sum_{x\in X_m}\pi(x)\log d_m(x)-\KL(\mu\|\mu^e).
\]
Equivalently,
\[
H(\mu)-H(\pi)
=
\mathbb{E}_{\pi}[\log d_m(X)]-\KL(\mu\|\mu^e).
\]
\end{proposition}

\begin{proof}
Write \(x_a:=\Fold_m(a)\). By Definition~\ref{def:fiber-uniform-lift},
\[
\mu^e(a)=\frac{\pi(x_a)}{d_m(x_a)}.
\]
Hence
\[
\begin{aligned}
\KL(\mu\|\mu^e)
&=
\sum_{a\in\Omega_m}\mu(a)\log\frac{\mu(a)}{\mu^e(a)}\\
&=
-H(\mu)-\sum_{a\in\Omega_m}\mu(a)\log\pi(x_a)
+\sum_{a\in\Omega_m}\mu(a)\log d_m(x_a).
\end{aligned}
\]
Since \(\sum_{a:\,x_a=x}\mu(a)=\pi(x)\), the two sums become
\[
\sum_{a\in\Omega_m}\mu(a)\log\pi(x_a)=\sum_{x\in X_m}\pi(x)\log\pi(x)=-H(\pi)
\]
and
\[
\sum_{a\in\Omega_m}\mu(a)\log d_m(x_a)
=
\sum_{x\in X_m}\pi(x)\log d_m(x).
\]
Rearranging gives the claim.
\end{proof}

\begin{proposition}[KL isometry of fiber-uniform lifts]\label{prop:fold-fiber-lift-kl-isometry}
For every \(\pi,\eta\in\Delta(X_m)\),
\[
\KL(Q_m\pi\|Q_m\eta)=\KL(\pi\|\eta).
\]
\end{proposition}

\begin{proof}
Write \(x_a:=\Fold_m(a)\). Then
\[
\frac{(Q_m\pi)(a)}{(Q_m\eta)(a)}
=
\frac{\pi(x_a)}{\eta(x_a)},
\]
so
\[
\begin{aligned}
\KL(Q_m\pi\|Q_m\eta)
&=
\sum_{a\in\Omega_m}\frac{\pi(x_a)}{d_m(x_a)}
\log\frac{\pi(x_a)}{\eta(x_a)}\\
&=
\sum_{x\in X_m}
\left(\sum_{a:\,x_a=x}\frac{\pi(x)}{d_m(x)}\right)
\log\frac{\pi(x)}{\eta(x)}
=
\KL(\pi\|\eta).
\end{aligned}
\]
\end{proof}

\begin{theorem}[Pythagorean decomposition for the canonical fold]\label{thm:fold-kl-pythagoras}
Let \(\mu\in\Delta(\Omega_m)\), let \(\pi=(\Fold_m)_*\mu\), and let
\(\mu^e=Q_m\pi\). Then for every \(\eta\in\Delta(X_m)\),
\[
\KL(\mu\|Q_m\eta)
=
\KL(\mu\|\mu^e)+\KL(\pi\|\eta).
\]
Consequently, \(\mu^e\) is the unique I-projection of \(\mu\) onto the
fiber-uniform model family
\[
\mathcal E_m:=\{Q_m\eta:\eta\in\Delta(X_m)\}\subset\Delta(\Omega_m),
\]
and
\[
\inf_{\eta\in\Delta(X_m)}\KL(\mu\|Q_m\eta)=\KL(\mu\|\mu^e).
\]
\end{theorem}

\begin{proof}
Let \(x_a:=\Fold_m(a)\). Since
\[
\mu^e(a)=\frac{\pi(x_a)}{d_m(x_a)},
\qquad
(Q_m\eta)(a)=\frac{\eta(x_a)}{d_m(x_a)},
\]
one has
\[
\log\frac{\mu^e(a)}{(Q_m\eta)(a)}=\log\frac{\pi(x_a)}{\eta(x_a)}.
\]
Therefore
\[
\begin{aligned}
\KL(\mu\|Q_m\eta)-\KL(\mu\|\mu^e)
&=
\sum_{a\in\Omega_m}\mu(a)\log\frac{\mu^e(a)}{(Q_m\eta)(a)}\\
&=
\sum_{a\in\Omega_m}\mu(a)\log\frac{\pi(x_a)}{\eta(x_a)}\\
&=
\sum_{x\in X_m}
\left(\sum_{a:\,x_a=x}\mu(a)\right)\log\frac{\pi(x)}{\eta(x)}
=
\KL(\pi\|\eta).
\end{aligned}
\]
This proves the decomposition. Since \(\KL(\pi\|\eta)\ge 0\), with equality if
and only if \(\eta=\pi\), the minimizer is unique and equals \(\mu^e=Q_m\pi\).
\end{proof}

\begin{proposition}[Dual characterization of the fiber-uniform lift]
\label{prop:dual-projection-maxent}
For any \(\pi\in\Delta(X_m)\), define the linear family
\[
\mathcal L_\pi:=\{\nu\in\Delta(\Omega_m):(\Fold_m)_*\nu=\pi\}.
\]
Then \(Q_m\pi\) is the unique maximum-entropy element of
\(\mathcal L_\pi\):
\[
H(\nu)\le H(Q_m\pi)
\qquad(\nu\in\mathcal L_\pi),
\]
with equality if and only if \(\nu=Q_m\pi\). Hence, for any
\(\mu\in\Delta(\Omega_m)\) with \(\pi=(\Fold_m)_*\mu\), the lift
\(Q_m\pi\) is simultaneously:
\begin{enumerate}
\item the maximum-entropy projection onto \(\mathcal L_\pi\), and
\item the I-projection of \(\mu\) onto the fiber-uniform family
  \(\mathcal E_m=\{Q_m\eta:\eta\in\Delta(X_m)\}\).
\end{enumerate}
\end{proposition}

\begin{proof}
For any \(\nu\in\mathcal L_\pi\),
Proposition~\ref{prop:fold-entropy-ledger} gives
\[
H(\nu)
=
H(\pi)+\sum_{x\in X_m}\pi(x)\log d_m(x)-\KL(\nu\|Q_m\pi).
\]
The first two terms are independent of \(\nu\), and
\(\KL(\nu\|Q_m\pi)\ge 0\), so \(H(\nu)\le H(Q_m\pi)\) with equality if and
only if \(\nu=Q_m\pi\). The I-projection statement is
Theorem~\ref{thm:fold-kl-pythagoras}.
\end{proof}

\begin{remark}[Relation to classical information projection]\label{rem:fold-kl-pythagoras-classical}
Theorem~\ref{thm:fold-kl-pythagoras} and
Proposition~\ref{prop:dual-projection-maxent} together realize the classical
duality of Csisz\'ar \cite{Csiszar1975IdivergenceGeometry} in the present
finite-state setting: the sufficient statistic is the folded observable
\(\Fold_m\), and the fiber-uniform family \(\mathcal E_m\) is exactly the
family of laws that are uniform within each fiber of this coarse-graining.
\end{remark}

\begin{corollary}[Microstate decomposition relative to the Parry lift]\label{cor:microstate-parry-lift-decomposition}
Let \(\mu\in\Delta(\Omega_m)\), let \(\pi=(\Fold_m)_*\mu\), and let
\(\mu^e=Q_m\pi\). Then
\[
\KL(\mu\|Q_m\pi_m)
=
\KL(\mu\|\mu^e)+\KL(\pi\|\pi_m).
\]
If \(|\supp(\pi)|\le k\), then
\[
\KL(\mu\|Q_m\pi_m)
\ge
\KL(\mu\|\mu^e)-\log M_m(k).
\]
In particular, assume \(\beta=\alpha\) and \(m\ge 6\). Then
\[
\KL(\mu_m\|Q_m\pi_m)
=
\mathbb E_{\pi_m^{\mathrm{fold}}}\big[\log d_m(Y)\big]
+\KL(\pi_m^{\mathrm{fold}}\|\pi_m),
\]
\[
\KL(\mu_m\|Q_m\pi_m)
\ge
\mathbb E_{\pi_m^{\mathrm{fold}}}\big[\log d_m(Y)\big]
+(m-1)\log\varphi-\log\pi(0)-\log(m+1),
\]
and for every \(N\ge 1\),
\[
\KL(\hat\mu_{m,N}\|Q_m\pi_m)
=
\mathbb E_{\hat\pi_{m,N}}\big[\log d_m(Y)\big]
+\KL(\hat\pi_{m,N}\|\pi_m),
\]
\[
\KL(\hat\mu_{m,N}\|Q_m\pi_m)
\ge
\mathbb E_{\hat\pi_{m,N}}\big[\log d_m(Y)\big]
+(m-1)\log\varphi-\log\pi(0)-\log(m+1).
\]
\end{corollary}

\begin{proof}
The first identity is Theorem~\ref{thm:fold-kl-pythagoras} with \(\eta=\pi_m\).
If \(|\supp(\pi)|\le k\), then
Theorem~\ref{thm:parry-support-barrier} gives
\(\KL(\pi\|\pi_m)\ge -\log M_m(k)\), which yields the second inequality.

Under \(\beta=\alpha\), Corollary~\ref{cor:sturmian-pure-degeneracy-ledger}
identifies \(\KL(\mu_m\|\mu_m^e)\) and \(\KL(\hat\mu_{m,N}\|\hat\mu_{m,N}^e)\)
with the corresponding degeneracy averages. The folded laws
\(\pi_m^{\mathrm{fold}}\) and \(\hat\pi_{m,N}\) are supported on at most \(m+1\)
words, and for \(m\ge 6\), one has \(m+1\le F_m\). Hence
Corollary~\ref{cor:exact-complexity-divergence-law} applies with \(k=m+1\),
yielding the lower bounds.
\end{proof}

\section{Local Zeckendorf Normalization and Uniqueness of the Fibonacci Weighting}\label{app:canonical-fold}

\subsection{Local Normalization Description of the Fold}

This subsection records a local description of the canonical fold from
Definition~\ref{def:canonical-fold}.

Let \(\mathcal{D}\) be the set of finitely supported digit configurations
\[
a=(a_1,a_2,\dots),\qquad a_k\in\mathbb{N},
\]
and define the Fibonacci value
\[
\mathrm{Val}(a):=\sum_{k\ge 1} a_k F_{k+1}.
\]
Embed \(\omega=\omega_1\cdots \omega_m\in\Omega_m\) as
\[
\iota_m(\omega):=(\omega_1,\dots,\omega_m,0,0,\dots)\in\mathcal{D}.
\]
Consider the following value-preserving local rewrites:
\[
e_k+e_{k+1}\to e_{k+2},
\qquad
2e_1\to e_2,
\qquad
2e_2\to e_1+e_3,
\qquad
2e_k\to e_{k-2}+e_{k+1}\ \ (k\ge 3),
\]
where \(e_k\) is the \(k\)-th unit vector in \(\mathcal{D}\).

\begin{proposition}[Order-independent local normalization]\label{prop:canonical-fold-rewrite}
Starting from \(\iota_m(\omega)\), every maximal rewrite sequence terminates at
the same normal form \(\mathrm{NF}(\iota_m(\omega))\in\{0,1\}^{\mathbb{N}}\)
with no adjacent \(1\)'s. This normal form is the Zeckendorf expansion
\(Z(N(\omega))\). Consequently,
\[
\Fold_m(\omega)=\pref_m\bigl(\mathrm{NF}(\iota_m(\omega))\bigr)
\]
is independent of rewrite order.
\end{proposition}

\begin{proof}
Each rewrite preserves \(\mathrm{Val}\) because it encodes one of the Fibonacci
identities
\[
F_{k+1}+F_{k+2}=F_{k+3},
\qquad
2F_2=F_3,
\qquad
2F_3=F_2+F_4,
\qquad
2F_{k+1}=F_{k-1}+F_{k+2}\ \ (k\ge 3).
\]
Termination follows from a lexicographic descent argument. Let
\[
A(a):=\sum_{k\ge 1} a_k,
\qquad
B(a):=\sum_{k\ge 1} k\,a_k,
\qquad
C(a):=a_2.
\]
The rules \(e_k+e_{k+1}\to e_{k+2}\) and \(2e_1\to e_2\) decrease \(A\); the
rules \(2e_k\to e_{k-2}+e_{k+1}\) for \(k\ge 3\) preserve \(A\) and decrease
\(B\); and the rule \(2e_2\to e_1+e_3\) preserves \(A\) and \(B\) but decreases
\(C\). Hence \((A,B,C)\) decreases lexicographically at each step, so rewriting
must terminate.

A normal form has no adjacent \(1\)'s and no digit \(\ge 2\), hence is a
Zeckendorf representation. Since every rewrite preserves \(\mathrm{Val}\) and
Zeckendorf representation is unique \cite{Zeckendorf1972}, all terminating
branches have the same normal form. The final formula for \(\Fold_m\) is then
just the definition of the canonical fold through prefix truncation.
\end{proof}

\subsection{Uniqueness of the Fibonacci Weighting}

The preceding subsection shows that the canonical fold is implemented by
Zeckendorf normalization. A natural question is whether the choice of
Fibonacci weights \(F_{k+1}\) is itself forced. The following theorem shows
that it is, within a natural axiom class.

\begin{theorem}[Contiguous numeration uniqueness]
\label{thm:fibonacci-weight-uniqueness}
Let \(w_2,w_3,\dots\) be positive integers. Define the valuation
\[
V_m^{(w)}(x_1\cdots x_m):=\sum_{k=1}^m x_k\,w_{k+1}.
\]
If for every \(m\ge 1\) the map \(V_m^{(w)}\) restricts to a bijection from
\(X_m\) onto \(\{0,1,\dots,F_{m+2}-1\}\), then \(w_{k+1}=F_{k+1}\) for all
\(k\ge 1\).
\end{theorem}

\begin{proof}
For \(m=1\), \(X_1=\{0,1\}\) and the bijection condition gives
\(\{0,w_2\}=\{0,1\}\), so \(w_2=1=F_2\).

For \(m=2\), \(X_2=\{00,01,10\}\) and \(\{0,w_2,w_3\}=\{0,1,2\}\); since
\(w_2=1\), one obtains \(w_3=2=F_3\).

Proceed by induction. Suppose \(w_2=F_2,\dots,w_m=F_m\) for some
\(m\ge 3\). Decompose \(X_m\) by its last symbols:
\[
X_m
=
\{u0: u\in X_{m-1}\}
\;\sqcup\;
\{v01: v\in X_{m-2}\}.
\]
On the first part, \(V_m^{(w)}(u0)=V_{m-1}^{(w)}(u)\). On the second,
\(V_m^{(w)}(v01)=V_{m-2}^{(w)}(v)+w_{m+1}\). By the inductive hypothesis,
\[
V_{m-1}^{(w)}(X_{m-1})=\{0,\dots,F_{m+1}-1\},
\quad
V_{m-2}^{(w)}(X_{m-2})=\{0,\dots,F_m-1\}.
\]
Hence
\[
V_m^{(w)}(X_m)
=
\{0,\dots,F_{m+1}-1\}
\;\sqcup\;
\{w_{m+1},\dots,w_{m+1}+F_m-1\}.
\]
For this to equal \(\{0,\dots,F_{m+2}-1\}=\{0,\dots,F_{m+1}+F_m-1\}\), the
second interval must start at \(F_{m+1}\), giving \(w_{m+1}=F_{m+1}\).
\end{proof}

\begin{corollary}[Uniqueness of the canonical residue fold]
\label{cor:canonical-fold-uniqueness}
Let \(G_m:\Omega_m\to X_m\) be a deterministic map satisfying:
\begin{enumerate}
\item \(G_m(x)=x\) for all \(x\in X_m\).
\item There exist positive integer weights \(w\) such that
  \(V_m^{(w)}|_{X_m}\) is a bijection onto \(\{0,\dots,F_{m+2}-1\}\).
\item For every \(\omega\in\Omega_m\),
  \(V_m^{(w)}(G_m(\omega))\equiv \sum_{k=1}^m\omega_k w_{k+1}
  \pmod{F_{m+2}}\).
\end{enumerate}
Then \(G_m=\Fold_m\).
\end{corollary}

\begin{proof}
Theorem~\ref{thm:fibonacci-weight-uniqueness} forces \(w_{k+1}=F_{k+1}\).
Condition~2 then provides a unique legal representative in each residue class
modulo \(F_{m+2}\), and condition~3 selects that representative. This is
exactly the residue characterization of
Proposition~\ref{prop:fold-residue-characterization}.
\end{proof}


\bibliographystyle{unsrt}
\bibliography{../../2026_golden_ratio_driven_scan_projection_generation_recursive_emergence/references,../../2026_golden_ratio_driven_scan_projection_generation_recursive_emergence/references_fibonacci,../../2026_golden_ratio_driven_scan_projection_generation_recursive_emergence/references_internal_2025}

\end{document}
